\documentclass[11pt]{article}
\usepackage[T1]{fontenc}
\usepackage{lmodern,microtype,amsmath,amssymb,amsthm,booktabs,array}
\usepackage[a4paper,margin=27mm]{geometry}
\usepackage[hidelinks]{hyperref}
\usepackage{longtable}
\setlength{\emergencystretch}{2em}
\newcommand{\workbenchurl}{https://github.com/KokunoYumeto/k-theory-and-galois-groups-workbench}
\newcommand{\recordurl}{https://github.com/KokunoYumeto/mathematics-workbenches}
\newcommand{\GT}{\widehat{\mathrm{GT}}}
\newcommand{\GTz}{\widehat{\mathrm{GT}}_0}
\newcommand{\F}{\mathbb F}
\newcommand{\Q}{\mathbb Q}
\newcommand{\Z}{\mathbb Z}
\newcommand{\grt}{\mathfrak{grt}_1}
\newcommand{\Lie}{\operatorname{Lie}}
\newcommand{\Hom}{\operatorname{Hom}}
\newcommand{\ad}{\operatorname{ad}}
\newcommand{\Ann}{\operatorname{Ann}}
\newcommand{\im}{\operatorname{im}}
\newcommand{\cts}{\mathrm{cts}}
\newtheorem{theorem}{Theorem}[section]
\newtheorem{proposition}[theorem]{Proposition}
\newtheorem{lemma}[theorem]{Lemma}
\newtheorem{corollary}[theorem]{Corollary}
\theoremstyle{definition}
\newtheorem{definition}[theorem]{Definition}
\theoremstyle{remark}
\newtheorem{remark}[theorem]{Remark}
\title{Abelian quotients separating the gentle and full\\Grothendieck--Teichm\"uller groups\\
\large After Ihara and Harbater--Schneps}
\author{}
\date{25 November 2026}
\begin{document}
\maketitle
\begin{abstract}
Drinfeld defined the Grothendieck--Teichm\"uller group by the cycle and pentagon relations, and Ihara's derivation algebra describes their infinitesimal counterparts. Harbater and Schneps identify the group defined by the two cycle relations with symmetric peripheral automorphisms; we use their description to construct its abelian quotients. We compute the indecomposable special derivations in
weight sixteen: their dimension is thirty-eight, over both the rationals
and the field with seventeen elements.  The raw count of possible brackets
overestimates their rank by one because of a Jacobi identity.  We also give
an exact weight-twelve pentagon certificate.  A precise
linear-algebra formulation identifies the functionals relevant to
characters: they annihilate the sum of the bracket space and the pentagon
kernel.  Using Brown's construction and the all-weight upper bound of OpenAI,
with Willwacher's subsequent proof,
and proving the required peripheral lifting and filtered group comparisons,
we obtain thirty-eight independent characters of order seventeen and
countably many independent characters of every finite $p$-power order
for each prime $p\geq5$, all vanishing on the full pentagon group.
A separate finite-cover construction gives countably many independent
quadratic characters with the same vanishing property.
We also extend the normalized elliptic-frame character to the whole
gentle group and construct an element on which it is nonzero while
the action on the third derived quotient is the identity. This character
is therefore independent of all the metabelian quadratic characters.
The mathematics
and text were produced by GPT-6 Astra Pro (ChatGPT Pro mode); the computations
were independently reimplemented and the proofs received separate checks
by the same model; no human has checked this paper.
\end{abstract}

\section{The question and the results}
Drinfeld defined the Grothendieck--Teichm\"uller group by compatibility
relations for braiding and reassociation \cite[\S5]{Drinfeld}.
Harbater and Schneps identify the group governed by the two cycle relations
with symmetric peripheral automorphisms of the four-point group
\cite[Main Theorem]{HS}.  Dolgushev and Guynee call it the gentle group;
their inverse-limit theorem distinguishes genuine elements from individual
finite approximations \cite[Definition 2.5, Theorem 5.2, Corollary 5.4]{DG}.
The full pentagon group is a closed subgroup of this gentle group.
The question here is what a continuous abelian character can detect when
the pentagon is omitted.

For precision, write $\widehat F_2=\langle X,Y\rangle$ and
$Z=(XY)^{-1}$.  An element of $\GTz$ is a pair
$(\lambda,f)\in\widehat\Z^\times\times\widehat F_2'$ for which
$X\mapsto X^\lambda$, $Y\mapsto f^{-1}Y^\lambda f$ is an automorphism
and, with $d=(\lambda-1)/2$,
\begin{align*}
 f(Y,X)f(X,Y)&=1,\\
 f(Z,X)Z^d f(Y,Z)Y^d f(X,Y)X^d&=1.
\end{align*}
By \cite[Theorems A and E]{Monoid}, every solution of these two
cycle equations with odd $\lambda$ already has unit multiplier and
defines an automorphism. Thus the definition is equivalently the
solution set of those equations. We use that theorem as cited
background, with its recorded AI checks and no human verification.
The subgroup $\GT$ additionally satisfies Drinfeld's pentagon in the
profinite pure braid group on four strands.  Its first homogeneous
relation is precisely \eqref{eq:pentagon} below.  The multiplier is
the homomorphism $(\lambda,f)\mapsto\lambda$.  The classical Galois
action gives $G_\Q\hookrightarrow\GT$ with cyclotomic multiplier;
in particular the multiplier on $\GT$ is onto
\cite[\S5]{Drinfeld}.

We use ordinary word weight: both generators of a free Lie algebra have
weight one.  Let $L=\Lie_k(x,y)$ and put $z=-x-y$.  For $n>1$, let $W_n$
consist of the homogeneous polynomials $q\in L_n$ satisfying
\begin{align}
 q(x,y)+q(y,x)&=0,\label{eq:anti}\\
 [x,q(x,y)]-[x+y,q(-x-y,y)]&=0.\label{eq:special}
\end{align}
Set $W_1=0$.  These are the normalized $S_3$-invariant outer special
derivations.  Their Lie bracket is the Ihara bracket
\begin{equation}
 \{q,r\}=[q,r]+D_q(r)-D_r(q),\qquad D_q(x)=0,\quad D_q(y)=[y,q].
 \label{eq:ihara}
\end{equation}
The cyclic identity $q(x,y)+q(y,z)+q(z,x)=0$ follows from
\eqref{eq:anti}--\eqref{eq:special}; it cannot replace
\eqref{eq:special} in this definition.

\begin{theorem}[Weight sixteen]\label{thm:weight16}
For $k=\Q$ and for $k=\F_{17}$,
\[
 \dim_k W_{16}=75,\qquad
 \dim_k[W,W]_{16}=37,\qquad
 \dim_k(W/[W,W])_{16}=38.
\]
The bracket space has a basis consisting of thirty-seven brackets of
homogeneous elements of weights less than sixteen.
\end{theorem}

Thus the special derivation equations leave thirty-eight independent
additive directions in this weight.  This statement concerns a graded Lie
algebra.  To obtain a character of a profinite group one must also realize
these directions in that group, pass through its filtered multiplication,
and establish invariance under the multiplier.  Section~\ref{sec:characters}
carries out these steps.

\begin{theorem}[Continuous character quotients]\label{thm:main-character-quotients}
Put $Q=\GTz/\overline{\langle\!\langle\GT\rangle\!\rangle}$.
There is a continuous surjection
\[
 Q^{\mathrm{ab}}\longrightarrow
 (\Z/2\Z)^{\mathbb N}\times
 \prod_{\substack{p\geq5\\p\ \mathrm{prime}}}\prod_{k,j\geq1}\Z/p^k\Z,
\]
already onto on the global multiplier-one subgroup. For $p=2$ and
for every prime $p\geq5$, the vector space
$\Hom_{\cts}(Q,\F_p)$ has dimension $\aleph_0$.
\end{theorem}

This theorem gives a large quotient of the abelianization; it does not
identify the entire quotient $Q$ or prove that $Q$ is abelian.
The quadratic factor is constructed in
Section~\ref{sec:quadratic-beta-characters} from finite-cover homology
and Ihara's beta factorization. Corollary~\ref{cor:simultaneous-quotient}
combines it with the odd-primary characters. Every character displayed
here vanishes on $\GT$, so it cannot distinguish $\GT$ from its
Galois subgroup.

\begin{theorem}[The elliptic-frame character]\label{thm:frame-overview}
There is a continuous character $\Delta:\GTz\to\F_2$ vanishing
on $\GT$ and an element $e=(1,f)\in\GTz$ such that
\[
 \rho_2(e)=0,\qquad f\in\widehat F_2''',\qquad \Delta(e)=1.
\]
Consequently $\Delta$ is outside the span of the quadratic characters
that factor through the metabelian action, including the beta characters
constructed in this paper.
\end{theorem}

Here $\rho_2$ is the integral Kummer cocycle defined by the pure
order-two frame in Section~\ref{sec:global-frame-character}.
That section fixes both frame signs and their arithmetic normalization,
then proves their invariance under every multiplier and the existence
of the global character. The proof of its vanishing on $\GT$ uses
the normalized vanishing theorem \cite[Theorem~3]{H17} as an input.
Section~\ref{sec:frame-witness} constructs the element $e$ using the
peripheral lifting theorem proved here. Its conjugator belongs to the
third derived subgroup, so its action on the full quotient
$\widehat F_2/\widehat F_2'''$ is the identity.

The input from the recent Deligne--Drinfeld theorem is generation of the
rational pentagon Lie algebra in odd weights.  OpenAI proved the all-weight
upper bound; Willwacher gives a shorter proof of that bound and explains
its combination with Brown's injectivity theorem
\cite[Theorem 1.1 and (7), pp.~1--2]{Willwacher}.
Markarian gives a geometric interpretation and a strict restricted tangent
reconstruction in the pro-two setting \cite[Theorem 4.8 and \S5]{Markarian}.
His result concerns that specified realization; the lifting argument below
addresses the full profinite group separately.
Our contribution is the peripheral lift, the transfer from even graded
indecomposables to full profinite characters, and the explicit weight-sixteen
calculation.  These characters quantify the gap created by omitting the
pentagon; they do not decide whether $G_{\Q}=\GT$.

Let $\mathfrak t_4$ be the infinitesimal pure braid Lie algebra, generated
by $t_{ij}=t_{ji}$ for $1\leq i<j\leq4$, with relations
\[
 [t_{ij},t_{kl}]=0\quad(i,j,k,l\text{ distinct}),\qquad
 [t_{ij},t_{ik}+t_{jk}]=0\quad(i,j,k\text{ distinct}).
\]
Define the linear pentagon map $P_n:W_n\to\mathfrak t_{4,n}$ by
\begin{align}
P_n(q)={}&q(t_{12},t_{23})+
q(t_{12}+t_{13},t_{24}+t_{34})+q(t_{23},t_{34})\notag\\
&-q(t_{12},t_{23}+t_{24})-q(t_{13}+t_{23},t_{34}).\label{eq:pentagon}
\end{align}

\begin{theorem}[The pentagon modulo seventeen]\label{thm:main-pentagon17}
Over $\F_{17}$ the map $P_{16}$ has rank seventy and a
five-dimensional kernel contained in $[W,W]_{16}$. Its restriction
to that bracket space has rank thirty-two. Consequently
\[
 W_{16}/[W,W]_{16}\longrightarrow
 \mathfrak t_{4,16}/P_{16}([W,W]_{16})
\]
is injective, of rank thirty-eight. Its dual injects linearly into
the continuous $\F_{17}$-valued characters of the full gentle group
which vanish on $\GT$, already after restriction to global
multiplier one. Every nonzero functional gives a character of order
seventeen.
\end{theorem}

Section~\ref{sec:mod17} proves both the exact finite-field ranks and
the integral passage from these functionals to group characters.
The rank assertion uses Brown's free motivic Lie algebra for
its upper bound and explicit modular minors for its lower bound.
The group construction uses the peripheral lifting theorem proved here.

\begin{proposition}[The character space of a weight]\label{prop:dual}
Put $B_n=[W,W]_n$ and $K_n=\ker P_n$.  The linear functionals on $W_n$
that annihilate both brackets and linear pentagon solutions are
\[
 C_n=\Ann(B_n)\cap\im(P_n^*)
     =\bigl(W_n/(B_n+K_n)\bigr)^*.
\]
In particular $P_n$ is a map on $W_n/B_n$ only when $P_n(B_n)=0$.
\end{proposition}

This distinction matters: a bracket of gentle directions need not satisfy
the pentagon.  One must therefore impose the pentagon kernel after passing
to indecomposables, or use the dual intersection in the proposition.

\begin{proposition}[An integral pentagon certificate]\label{prop:weight12}
The space $W_{12}$ has dimension seven and its bracket subspace has
dimension six.  For the integral basis specified in Section~\ref{sec:cert},
the primitive functional
\[
\ell=(-4,1,8,-4,-17,14,84)
\]
annihilates both $[W,W]_{12}$ and $\ker P_{12}$ over $\Q$.
It has the same properties after reduction modulo thirteen and is nonzero
there.  Consequently $C_{12}$ is one-dimensional over $\Q$ and $\F_{13}$.
\end{proposition}

\begin{theorem}[Integral transfer at top weight]\label{thm:main-topweight}
For every prime $p\geq5$, set $m=p-1$ and form $W_m$, $B_m=[W,W]_m$
and $K_m=\ker P_m$ over $\F_p$. After a choice of tame splitting there
is an injective linear map
\[
 (W_m/(B_m+K_m))^*
 \longrightarrow \Hom_{\cts}(Q,\F_p).
\]
The character associated with $\ell$ restricts to $\ell$ on the
leading weight-$m$ subgroup, using its identification with $W_m$.
These characters are independent already on the global multiplier-one
subgroup. In particular there is a nonzero character of order thirteen
from weight twelve, and thirty-eight independent characters of order
seventeen from weight sixteen.
\end{theorem}

The tame splitting separates the nonzero weight characters of the
cyclic multiplier group $\F_p^\times$ in the finite jet; its construction
and the leading-weight identification are given in
Theorem~\ref{thm:topweight}. This transfer theorem and its two stated
finite consequences do not use the Deligne--Drinfeld generation theorem.

The lifting result used throughout has a relative form. In the free case
let $\Lambda=\widehat F_2=\langle x_1,x_2\rangle$; in the modular case
let $\Lambda=\widehat{C_2*C_3}=\langle x_1,x_2\mid x_1^2=x_2^3=1\rangle$.
In either case put $x_3=(x_1x_2)^{-1}$. For a closed normal subgroup
$U\triangleleft\Lambda$, write $\gamma_n(U)$ for its closed lower
central series, put $R=\bigcap_{n\geq1}\gamma_n(U)$, and let
$\pi:\Lambda\to\overline\Lambda=\Lambda/R$ be the quotient map.

\begin{theorem}[Peripheral lifting with fixed conjugators]
\label{thm:main-peripheral-lifting}
Let $\alpha$ be a continuous automorphism of $\overline\Lambda$ and
let $\overline c_i\in\overline\Lambda$ satisfy
\[
 \alpha(\pi(x_i))=\overline c_i\pi(x_i)\overline c_i^{-1}
 \qquad(1\leq i\leq3).
\]
There are a continuous automorphism $\Phi$ of $\Lambda$ and
$c_i\in\Lambda$ such that
\[
 \pi\Phi=\alpha\pi,\qquad \pi(c_i)=\overline c_i,\qquad
 \Phi(x_i)=c_ix_ic_i^{-1}\quad(1\leq i\leq3).
\]
If one specified $\overline c_i$ is the identity, the lift can be
chosen with that $c_i$ equal to the identity.
\end{theorem}

Section~\ref{sec:peripheral-lifting} proves this theorem for every
closed normal $U$ in either source group. It provides an automorphism
of the full profinite group with the prescribed quotient action and
conjugator images, which is the step needed to realize the characters.

Furusho's theorem proves that the pentagon implies the hexagons in the
Lie, unipotent, pro-prime and pronilpotent settings
\cite[Theorems 1--3 and the two corollaries in \S3]{Furusho}.
It does not assert the converse.
Our explicit polynomial in Section~\ref{sec:seven} exhibits the failure
of the converse for invariant special derivations.  Brown's theorem
constructs the free motivic Lie algebra inside the associator symmetry
algebra \cite[Theorem 1.1 and its consequence, \S1.1]{Brown}; it does not identify an arbitrary gentle
direction with a motivic one.  These distinctions are necessary when
transferring a calculation to the full profinite group.

\section{Proof overview}
There are two separate algebraic questions.  First, a normalized special
derivation is determined by one Lie polynomial, and invariance under the
permutations of three punctures imposes explicit linear equations.  The
Witt character formula computes the dimension of their solution space.
Second, the abelianization requires the rank of the Ihara brackets.  In
weight sixteen the possible pairs give thirty-eight columns, but Jacobi
forces one relation.  A minor that is nonzero modulo seventeen proves
that this is the only relation, in characteristic zero as well.

For the pentagon we use its action on the free Lie algebra on three
letters.  A coefficient of a composition of derivations can be computed
by applying half of the operators forward and contracting the other half
backward.  This gives exact integer rows without expanding every output
polynomial.  The weight-twelve functional is an explicit integer linear
combination of five such rows.  The argument establishes vanishing on
the whole pentagon kernel, rather than only on selected examples.

A finite jet is not yet a full profinite automorphism. The lifting
argument compares peripheral homomorphism counts in every subgroup of
a finite quotient. Finite central extensions encode the irreducible
representation data that enter those counts. Equality of all subgroup
counts preserves the positive number of epimorphisms, and compactness
then gives a single automorphism satisfying every finite requirement.
This proves the peripheral lifting theorem used in both constructions.

For odd-primary characters, the local multiplier-one image is a compact
unipotent group in each bounded weight. A Galois element with multiplier
$1+p$ splits its rational Lie algebra into distinct weight eigenspaces.
Even indecomposable functionals vanish on the pentagon Lie algebra by
the Deligne--Drinfeld theorem. Their compact additive images give integral
lattices, and the condition $(p-1)p^{k-1}\mid m$ makes their reductions
invariant under every multiplier. Arithmetic normalization then extends
them to the full gentle group. At top weight $p-1$, tame integral
projectors give the more precise construction for each modular functional.

At two, the degree-thirty-six cover has a distinguished pair of cubic
character lines. Changing the sign on just that pair acts trivially
on the finite quotient's first and second integral homology. The sign
therefore lifts through every finite central extension with the required
peripheral conjugators. Its beta-associativity defect has a norm-one
quadratic sign, yielding the independent family. Finally, a different
relative peripheral lift places a central modular frame value in the
third derived subgroup. It detects the global frame character while
all metabelian characters vanish.

\section{Special derivations and their dimensions}\label{sec:dimensions}
This section identifies the equations used in every calculation.
An outer special derivation is the class, modulo inner derivations, of a
derivation tangent to each of $x,y,z$.  For degree $n>1$, subtracting an
inner derivation gives a unique representative $D_q(x)=0$,
$D_q(y)=[y,q]$.  Uniqueness follows because the centralizer of a free
generator has no homogeneous Lie component of degree greater than one.
For completeness, in the free associative algebra the leading and trailing
letters in $ux=xu$ force $u$ to be a polynomial in $x$; the free Lie algebra
on one generator has only its degree-one part.  This argument is valid
over $\Z$ and after reduction modulo any prime.

Let $s$ exchange $x,y$, and let $r$ send $(x,y,z)$ to $(y,z,x)$.
Directly evaluating the difference $sD_qs^{-1}-D_q$ on $x,y$ shows that
it is inner precisely when \eqref{eq:anti} holds.  Evaluating
$rD_qr^{-1}-D_q$ gives the additional condition
\[
 [x,q(x,y)]=[z,q(y,z)],
\]
which is \eqref{eq:special} after antisymmetry.  The three cyclic
transforms imply that the adjoint action of
$q(x,y)+q(y,z)+q(z,x)$ is zero.  The free Lie algebra has zero centre,
so the cyclic identity follows.  The same identities show tangency at
$z$, and prove the converse.  Commutators of special derivations remain
special, and invariant outer classes remain invariant; evaluation on $y$
gives \eqref{eq:ihara}.

Write $V$ for the standard two-dimensional representation of $S_3$, with
coordinates $x+y+z=0$.  The permutation representation on the three
peripheral conjugators modulo diagonal inner conjugators is $V$.
The bracket map is onto, because repeated Jacobi identities give
$[L_n,x]+[L_n,y]=L_{n+1}$.  Therefore there is an exact sequence
\begin{equation}
0\longrightarrow\operatorname{OutSDer}_n
 \longrightarrow V\otimes L_n(V)\longrightarrow L_{n+1}(V)
 \longrightarrow0.\label{eq:exact}
\end{equation}
It is exact over $\Z$; for primes different from two and three, taking
$S_3$ invariants is exact by averaging.  Thus these kernels commute with
extension from $\Z_p$ to $\Q_p$ and reduction to $\F_p$.

For $g\in S_3$, put
\[
 l_n(g)=\frac1n\sum_{d\mid n}\mu(d)
                  \operatorname{tr}(g^d\mid V)^{n/d}.
\]
The Witt character identity follows by applying logarithms to the
Poincar\'e--Birkhoff--Witt product for the tensor algebra and then using
M\"obius inversion.  It gives $l_n(g)=\operatorname{tr}(g\mid L_n(V))$.
For a transposition $s$ and three-cycle $r$, the traces on $V$ are $0,-1$,
respectively.  Taking invariants in \eqref{eq:exact} consequently gives
\begin{equation}
 a_n:=\dim W_n=
 \frac{2l_n(1)-2l_n(r)-l_{n+1}(1)-3l_{n+1}(s)-2l_{n+1}(r)}6.
 \label{eq:dimension}
\end{equation}
Here and below the expression is first evaluated as an integer in
characteristic zero; the integral exact sequence proves the same
dimensions in characteristic $p\geq5$.

\begin{center}
\begin{tabular}{c|rrrrrrrrrrrrrr}
$n$&3&4&5&6&7&8&9&10&11&12&13&14&15&16\\\hline
$a_n$&1&0&1&0&2&1&4&2&9&7&21&24&56&75
\end{tabular}
\end{center}
The elementary upper bound on the bracket rank is
\begin{equation}
 b_n=\sum_{2\leq i<n-i}a_i a_{n-i}
       +\begin{cases}\binom{a_{n/2}}2,&2\mid n,\\0,&2\nmid n.
       \end{cases}\label{eq:raw}
\end{equation}
This counts possible pairs; it does not assert their independence.

Antisymmetry by itself is not the standard-isotypic condition.  For
example $[x,y]$ is antisymmetric, but its cyclic sum is $3[x,y]$.
Even the two conditions together do not imply specialness.  The
polynomial
\[
 q=-b_{xxxy}+4b_{xxyy}-b_{xyyy}
\]
has $sq=-q$ and $(1+r+r^2)q=0$, while the coefficient of $x^4y$ in
$[x,q(x,y)]-[x+y,q(-x-y,y)]$ is $-2$.  This is an explicit
weight-four distinction over $\Q$ and in every characteristic at least
five.  The bracket-closed space used throughout this paper is $W$,
with the special equation imposed.

\section{The thirty-seven brackets in weight sixteen}\label{sec:brackets}
This section proves Theorem~\ref{thm:weight16}, including the dependence
missed by the raw pair count.
Only the pairs $(3,13)$, $(5,11)$ and $(7,9)$ can contribute: their
dimensions are $21,9,8$.  The pair $(8,8)$ contributes zero because
$W_8$ has dimension one.  Thus there are thirty-eight columns.

Let $b_w$ denote the standard Lyndon bracket, for $x<y$, using the longest
proper Lyndon suffix at every step.  Set
\begin{align*}
 u&=-b_{xxy}+b_{xyy},\\
 v&=-2b_{xxxxy}+4b_{xxxyy}-3b_{xxyxy}
       -4b_{xxyyy}-b_{xyxyy}+2b_{xyyyy}.
\end{align*}
Both satisfy \eqref{eq:anti}--\eqref{eq:special}.  Jacobi gives
\begin{equation}
 \{u,\{v,\{u,v\}\}\}=\{v,\{u,\{u,v\}\}\}.
 \label{eq:jacobi}
\end{equation}
The inner terms are nonzero.  Their respective leading associative words
and coefficients are
\[
 \{u,v\}:\ -4x^5yxy,\quad
 \{u,\{u,v\}\}:\ 4x^7yxyy,\quad
 \{v,\{u,v\}\}:\ 8x^9yxyy.
\]
In particular the relation \eqref{eq:jacobi} is a nonzero relation
between the $(3,13)$ and $(5,11)$ columns, over $\Q$ and $\F_{17}$.
The bracket rank is at most thirty-seven.

Here are complete specifications for the finite lower-bound calculation.
List Lyndon words in lexicographic order.  For every $n\leq13$, form a
matrix whose column at $b_w$ is the list of Lyndon-word coefficients of
$b_w+b_w(y,x)$ followed by those of
$[x,b_w]-[x+y,b_w(-x-y,y)]$.
The nullspace in reduced row echelon form, with free columns in increasing
order, gives a rational basis.  Clear denominators in each vector and
divide by the greatest common divisor of its entries.
Each vector is checked against every original row.  A complementary-rank
minor modulo seventeen proves equality with the full rational kernel.
The resulting bases have the dimensions in \eqref{eq:dimension}, and
their reductions remain independent modulo both thirteen and seventeen.

Use these bases, in order, for the thirty-eight bracket columns described
above.  Evaluate each bracket at every length-sixteen Lyndon word.  The
columns numbered $1,\ldots,38$ except column $28$ are independent: the
lexicographically first pivot rows modulo seventeen give a square minor
whose determinant is $3$ modulo seventeen.  The pivot words are specified
in the certificate by the following lengths-sixteen binary integers,
where $x=0$, $y=1$:
\begin{center}
\begin{minipage}{0.93\linewidth}\small\ttfamily\raggedright
5, 17, 23, 27, 39, 43, 71, 75, 77, 83, 85, 95, 99, 111, 139, 141, 147, 149, 159, 163, 175, 183, 187, 215, 275, 303, 311, 315, 317, 343, 347, 349, 359, 365, 599, 619, 683.

\end{minipage}
\end{center}
The matrix entries are integers determined by the preceding definitions,
so this nonzero minor proves rank at least thirty-seven in both fields.
Together with \eqref{eq:jacobi} and $a_{16}=75$, it proves the theorem.

The calculation uses only the following elementary triangular fact.
The associative expansion of $b_w$ has coefficient one at $w$ and has no
word smaller than $w$.  Hence the restriction to Lyndon-word coefficients
is an invertible triangular map on $L_n$.  A Lie polynomial is zero if
and only if all those coefficients are zero.  The matrix construction
therefore omits no defining equation.

\section{Exact pentagon certificates}\label{sec:cert}
This section proves Proposition~\ref{prop:weight12} and describes the
finite calculation independently of a particular program.
The integral seven-element basis of $W_{12}$ is obtained by precisely the
reduced-row-echelon prescription in Section~\ref{sec:brackets}.  The six
brackets have rank six, with a minor nonzero modulo thirteen, and their
coordinate columns are annihilated by $\ell$.

Let $a,b,c$ freely generate a Lie algebra.  For $i,j\leq3$, let $t_{ij}$
send its $i$th generator to its bracket with the $j$th, the $j$th generator
to its bracket with the $i$th, and the remaining generator to zero.
Let $t_{14}=\ad(a)$, $t_{24}=\ad(b)$, $t_{34}=\ad(c)$.
Evaluation on the three generators verifies the infinitesimal braid
relations, so this defines a representation of $\mathfrak t_4$.
Let $R_1,\ldots,R_5$ be the following coefficients of $P_{12}(q)$ acting
in this representation:
\begin{center}
\begin{tabular}{ccl}
\toprule
row&input generator&output word\\\midrule
1&$a$&\texttt{aaaaaaabcbcbc}\\
2&$a$&\texttt{acbaacaaabcba}\\
3&$b$&\texttt{aaaabacaacabb}\\
4&$b$&\texttt{baaacacbaabac}\\
5&$a$&\texttt{abbacbaaacaca}\\\bottomrule
\end{tabular}
\end{center}
Their rows in the prescribed basis are
\[
\begin{pmatrix}
180&-110&-160&296&214&-364&-167700\\
-13609&3331&-582&8438&-6239&-1104&-900331\\
-2160&1200&840&-960&-1800&2400&1150800\\
480&-300&-360&540&540&-840&-394680\\
-2702&1870&3472&-6964&-3554&7240&3662710
\end{pmatrix}.
\]
Ordinary multiplication verifies
\begin{equation}
698922\ell=-192329R_1+472R_2-26842R_3-107926R_4-11870R_5.
\label{eq:certificate12}
\end{equation}
The integer $698922$ is a unit modulo thirteen.  This proves that $\ell$
annihilates the pentagon kernel over both fields.  The bracket calculation
proves that its nonzero line is exactly $\Ann([W,W]_{12})$, so the
claimed dimension of $C_{12}$ follows.

The coefficient rows can be checked by a short exact recurrence.
If a derivation $D$ replaces a letter by a two-letter polynomial, then
\begin{equation}
[w]D(A)=\sum_{j,v}[w_jw_{j+1}]D(v)
 [w_1\cdots w_{j-1}v w_{j+2}\cdots w_m]A.
\label{eq:contraction}
\end{equation}
This is the Leibniz rule with terms grouped by the adjacent pair that
was created.  Split each twelve-letter operator word into two blocks of
six.  Apply the second block to the input generator by the Leibniz rule;
apply the transpose of the first block to the chosen output word by
\eqref{eq:contraction}; then sum the products of matching middle-word
coefficients.  The result is exactly its coefficient in the composition.
Summing with the associative coefficients of $q$ and the five signs in
\eqref{eq:pentagon} produces the displayed integer matrix.
All intermediate operations are integral.  In particular, searching for
rows modulo a prime is not used to justify an identity over $\Q$.

\section{A degree-seven witness}\label{sec:seven}
This section gives a small exact example showing why the pentagon imposes
an additional condition.  Put
\begin{align*}
h={}&-3b_{xxxyxyy}-3b_{xxxyyxy}-b_{xxyxyxy}
       +3b_{xxyxyyy}\\
 &+3b_{xxyyxyy}+3b_{xxyyyxy}+b_{xyxyxyy}.
\end{align*}
Expanding the seven standard brackets gives sixty-two nonzero associative
coefficients.  Substitution gives
\[
h(x,y)+h(y,x)=0,\qquad
[x,h(x,y)]=[x+y,h(-x-y,y)],
\]
and hence $h\in W_7$.  Consider the map from $\mathfrak t_4$ to the
commutator Lie algebra of $\Z[S_4]$ given by $t_{ij}\mapsto(ij)$.
Disjoint transpositions commute, and
$[(ij),(ik)+(jk)]=0$ by multiplication of the three transpositions.
The map is therefore a Lie homomorphism.  Substitution gives
\begin{equation}
 P_7(h)=-160\bigl((34)-(12)\bigr)\ne0.
 \label{eq:seven}
\end{equation}
This identity is checked in the twenty-four-element group-ring basis;
the two displayed transpositions are its only nonzero coefficients.
It remains nonzero modulo eleven, as does twice the identity.
Thus $W_7$ strictly contains the linear pentagon solution space.
The passage from this polynomial to a full profinite pair requires the
peripheral lifting theorem; the finite polynomial calculation alone does
not make that passage.

\section{Which linear functionals can give characters}\label{sec:linear}
This section proves Proposition~\ref{prop:dual} without imposing an
unnecessary descent condition on the pentagon map.
For a linear map $P:E\to T$ of finite-dimensional vector spaces,
$\im P^*=\Ann(\ker P)$: the forward inclusion is immediate, and a
functional vanishing on the kernel factors through $\im P$ and extends
linearly to $T$.  Intersecting with $\Ann(B)$ gives
$\Ann(B+\ker P)$, which is the dual of $E/(B+\ker P)$.
Finally, $P$ factors through $E/B$ precisely when it vanishes on $B$.
This proves every assertion of the proposition.

For the group-theoretic use, take a prime $p\geq5$, and suppose the
multiplier-one image of $\GTz$ in pro-$p$ automorphisms has a filtration
$U_n$ with
\[
[U_i,U_j]\subset U_{i+j},\qquad U_n/U_{n+1}\simeq W_n(\Z_p).
\]
A functional on a single associated-graded piece need not by itself be a
character of $U$.  The following sections establish the displayed
filtration, split the filtered Lie algebra, normalize its additive image,
and extend the resulting homomorphisms through the multiplier group.

\IfFileExists{lifting.tex}{\section{Peripheral automorphisms and the pronilpotent quotient}
\label{sec:peripheral-lifting}

We prove a lifting theorem with prescribed peripheral conjugators.
Keeping these conjugators is necessary when a peripheral element must
remain fixed after lifting. The finite argument uses central extensions
to retain the character values which control the relevant lifting fibres.

There are two source groups. In the free case set
\[
 \Lambda=\widehat F_2=\langle x_1,x_2\rangle,
 \qquad x_3=(x_1x_2)^{-1},\qquad U=\Lambda.
\]
In the modular case set
\[
 \Lambda=\widehat{C_2*C_3}
     =\langle x_1,x_2\mid x_1^2=x_2^3=1\rangle,
 \qquad x_3=(x_1x_2)^{-1},
\]
and let $U$ be the kernel of the standard quotient $\Lambda\to S_3$.
The theorem below also holds for any closed normal subgroup $U$ of
either source; these two choices are its initial applications.
For such a subgroup, let $\gamma_n(U)$ be its closed lower central series,
and put
\[
 R=\bigcap_{n\geq1}\gamma_n(U),\qquad
 \overline\Lambda=\Lambda/R,
 \qquad \pi:\Lambda\longrightarrow\overline\Lambda.
\]
A positive peripheral automorphism sends each $x_i$ to a conjugate of
$x_i$, with no change of exponent.

\begin{theorem}[Peripheral lifting]\label{thm:peripheral-lifting}
Let $\alpha$ be a continuous automorphism of $\overline\Lambda$.
Suppose that elements $\overline c_i\in\overline\Lambda$ satisfy
\[
 \alpha(\pi(x_i))
    =\overline c_i\pi(x_i)\overline c_i^{-1}
       \qquad(1\leq i\leq3).
\]
There are a continuous automorphism $\Phi$ of $\Lambda$ and elements
$c_i\in\Lambda$ such that
\[
 \pi\Phi=\alpha\pi,
 \qquad \pi(c_i)=\overline c_i,
 \qquad \Phi(x_i)=c_ix_ic_i^{-1}
       \qquad(1\leq i\leq3).
\]
If one specified conjugator $\overline c_i$ is the identity, the lift
can be chosen so that the corresponding $c_i$ is the identity.
\end{theorem}

The last assertion normalizes one chosen conjugator. It does not claim
that several full conjugators can be normalized independently by a
single inner automorphism.

\subsection{A finite central refinement}

\begin{lemma}\label{lem:finite-peripheral-refinement}
Let $q:\Lambda\twoheadrightarrow E$ be a finite quotient, let
$A\triangleleft E$, and write $D=E/A$ and $q_D=q\bmod A$.
There are a finite central extension
$\beta:D^\dagger\twoheadrightarrow D$ and an epimorphism
$q^\dagger:\Lambda\twoheadrightarrow D^\dagger$ such that
$\beta q^\dagger=q_D$, with the following property.

Suppose $q_1^\dagger:\Lambda\twoheadrightarrow D^\dagger$ is an
epimorphism and $r_i\in D^\dagger$ satisfy
\[
 q_1^\dagger(x_i)=r_iq^\dagger(x_i)r_i^{-1}
       \qquad(1\leq i\leq3).
\]
There is an epimorphism $q_1:\Lambda\twoheadrightarrow E$ and
elements $C_i\in E$ with
\[
 q_1(x_i)=C_iq(x_i)C_i^{-1},\qquad
 C_i\bmod A=\beta(r_i),\qquad
 q_1\bmod A=\beta q_1^\dagger.
\]
\end{lemma}

\begin{proof}
Put $X_i=q(x_i)$ and let
\[
 T_i=\{aX_ia^{-1}:a\in A\}
\]
be its $A$-conjugacy class. We shall retain enough information to
compare homomorphism counts into every subgroup of $E$, and then
recover the number of epimorphisms.

Fix a subgroup $J\leq E$ mapping onto $D$ such that each $T_i\cap J$
is nonempty, and put $M=J\cap A$. Let
$\rho:M\to\operatorname{GL}(V_\rho)$ be an irreducible complex
representation whose isomorphism class is fixed by conjugation by $J$.
Write $d_\rho=\dim V_\rho$ and $e=\exp J$. Define
\[
 \widetilde J_\rho=
 \left\{(j,B)\in J\times\operatorname{GL}(V_\rho):
 \begin{array}{l}
 B\rho(m)B^{-1}=\rho(jmj^{-1})\quad(m\in M),\\
 \det B\in\mu_e
 \end{array}\right\}.
\]
The invariance of $\rho$ gives an implementing matrix above every $j$.
By Schur's lemma all such matrices form one scalar coset. In that
coset the determinant condition has exactly $d_\rho e$ solutions:
if $B_0$ is one implementing matrix, it is the equation
$t^{d_\rho e}=(\det B_0)^{-e}$ for $B=tB_0$. Matrix multiplication
and inversion preserve the defining conditions. Thus
$\widetilde J_\rho$ is a finite group, and its kernel over $J$ is
the central scalar group $\mu_{d_\rho e}$.

The map
\[
 M\longrightarrow\widetilde J_\rho,
 \qquad m\longmapsto(m,\rho(m))
\]
is a normal embedding. Its matrices satisfy the determinant condition
because $m^e=1$, and normality is exactly the intertwining identity.
This copy of $M$ intersects the scalar kernel trivially. Consequently
\[
 Q_\rho=\widetilde J_\rho/M
\]
is a finite central extension of $J/M=D$. The natural map gives an
isomorphism
\begin{equation}
 \widetilde J_\rho\simeq J\times_DQ_\rho.
 \label{eq:representation-fibre-product}
\end{equation}
Indeed it is injective, and any pair with matching images in $D$ is
obtained by multiplying a representative in $\widetilde J_\rho$ by
an element of the embedded $M$. The matrix coordinate therefore gives
a representation
\[
 \Psi_\rho:J\times_DQ_\rho
        \longrightarrow\operatorname{GL}(V_\rho)
\]
whose restriction to $M\times\{1\}$ is $\rho$.

We next choose a source homomorphism
$R_\rho^0:\Lambda\to Q_\rho$ lifting $q_D$. For the free source,
choose arbitrary lifts of $q_D(x_1)$ and $q_D(x_2)$. For the modular
source choose $j_i\in J\cap T_i$ for $i=1,2$. The orders of $j_1$
and $j_2$ divide two and three, respectively. If $j_i$ has order
$h>1$, an implementing matrix has its $h$th power scalar by Schur's
lemma. Rescaling makes that power the identity. Its determinant then
has order dividing $h$, and hence lies in $\mu_e$. If $j_i=1$, use
the identity matrix. The resulting elements of $Q_\rho$ have orders
dividing two and three, respectively, and define the required
homomorphism from the modular source.

Perform this construction for all the finitely many relevant $J$ and
all their invariant irreducible representations. Take the joint image
of the maps $q_D$ and $R_\rho^0$ in their fibre product over $D$.
Call it $D^\dagger$, and call the resulting source epimorphism
$q^\dagger$. Its kernel over $D$ is a subgroup of a product of central
scalar groups, so $D^\dagger\to D$ is a finite central extension.

We now establish the counting identity used to transport this data.
Fix one of the subgroups $J$ above. Let $\chi:\Lambda\to D$ be
an epimorphism. Suppose $V_i\subset J$ is invariant under conjugation
by $M$ and lies in the fibre over $\chi(x_i)$. For each invariant
$\rho$ suppose $R_\rho:\Lambda\to Q_\rho$ is a homomorphism lifting
$\chi$. Then the number of triples with product one is
\begin{equation}
 \begin{split}
 &\#\{(a_1,a_2,a_3)\in V_1\times V_2\times V_3:
                   a_1a_2a_3=1\}\\
 &\qquad=\frac1{|M|}
  \sum_{\rho\in\operatorname{Irr}(M)^J}\frac1{d_\rho}
  \prod_{i=1}^3
    \left(\sum_{a\in V_i}
       \operatorname{tr}\Psi_\rho(a,R_\rho(x_i))\right).
 \end{split}
 \label{eq:relative-peripheral-count}
\end{equation}
To prove this, put $B_i=\sum_{a\in V_i}a$ in $\mathbb C[J]$.
Each $B_i$ commutes with $\mathbb C[M]$. Let $e_\tau$ be the
primitive central idempotent of $\mathbb C[M]$ belonging to an
irreducible representation $\tau$. Since every term of $B_i$ has
the same image in $D$, multiplication by $B_i$ takes the $\tau$
block to the block obtained by conjugating $\tau$ by that image.
But $e_\tau B_i=e_\tau B_ie_\tau$, by commutation. Hence this
product is zero unless $\chi(x_i)$ fixes $\tau$. The images of
$x_1,x_2$ generate $D$, so only $J$-invariant representations can
contribute to $B_1B_2B_3$.

This product lies in $\mathbb C[M]$, because the three quotient
values multiply to one. The coefficient of the identity in an
element $B\in\mathbb C[M]$ is
\[
 \frac1{|M|}\sum_{\tau\in\operatorname{Irr}(M)}
                  (\dim\tau)\operatorname{tr}\tau(B),
\]
by the regular representation of $M$. For an invariant $\rho$, set
\[
 S_{i,\rho}=\sum_{a\in V_i}\Psi_\rho(a,R_\rho(x_i)).
\]
Conjugating by $\rho(m)$ permutes this sum, so Schur's lemma makes
it a scalar matrix, namely
$d_\rho^{-1}\operatorname{tr}(S_{i,\rho})I$.
Since $R_\rho(x_1)R_\rho(x_2)R_\rho(x_3)=1$, multiplication in
the fibre product gives
\[
 \rho(B_1B_2B_3)=S_{1,\rho}S_{2,\rho}S_{3,\rho}.
\]
The three scalar denominators and the two dimension factors in the
regular-character calculation give the factor $1/d_\rho$ in
\eqref{eq:relative-peripheral-count}. This proves the identity.
For the modular source, when $V_1,V_2$ are subsets of the prescribed
conjugacy classes, their elements already satisfy the square and cube
relations. In that case the same triple count is the number of source
homomorphisms with these conditions.

Return now to $q_1^\dagger$ and its specified conjugators $r_i$.
Choose $H_i\in E$ mapping to $\beta(r_i)$ and put
$T_i^1=H_iT_iH_i^{-1}$. For every $J\leq E$ mapping onto $D$,
choose $j_i\in J$ with the same image as $H_i$ in $D$.
Writing $H_i=j_ia_i$ with $a_i\in A$ gives the exact equality
\begin{equation}
 T_i^1\cap J=j_i(T_i\cap J)j_i^{-1}.
 \label{eq:relative-class-intersection}
\end{equation}
Here we use that $T_i$ is an entire $A$-conjugacy class; conjugation
by $a_i$ preserves it, and conjugation by $j_i$ preserves $J$.

For a relevant $J$ and $\rho$, let $r_{i,\rho}$ be the projection
of $r_i$ to $Q_\rho$, and let $R_\rho^1$ be the corresponding
projection of $q_1^\dagger$. The pair $(j_i,r_{i,\rho})$ belongs
to $J\times_DQ_\rho$. It conjugates the old pairs
\[
 (a,R_\rho^0(x_i)),\quad a\in T_i\cap J,
\]
to the new pairs
\[
 (j_iaj_i^{-1},R_\rho^1(x_i)),\quad
       j_iaj_i^{-1}\in T_i^1\cap J.
\]
The trace sums in \eqref{eq:relative-peripheral-count} are therefore
equal term by term. The number of homomorphisms $\Lambda\to J$
with reduction $q_D$ and values in $T_i$ equals the number with
reduction $\beta q_1^\dagger$ and values in $T_i^1$.
If some $T_i\cap J$ is empty, both counts are zero by
\eqref{eq:relative-class-intersection}. If $J$ does not map onto $D$,
both counts are zero because both specified quotient maps are onto.
Thus the equality holds for every subgroup $J\leq E$.

Each homomorphism into $J$ has a unique image subgroup. M\"obius
inversion on the finite subgroup lattice, or induction on subgroup
order, recovers the number of epimorphisms to $E$ from these counts.
The two epimorphism counts agree. The old count is positive because
it contains $q$. Hence a new epimorphism $q_1$ exists. For each $i$,
membership $q_1(x_i)\in T_i^1$ gives an element $a_i'\in A$ such
that
\[
 q_1(x_i)=H_ia_i'X_i(a_i')^{-1}H_i^{-1}.
\]
Taking $C_i=H_ia_i'$ gives the prescribed conjugator image in $D$.
The quotient identity follows on the two generators, proving the lemma.
\end{proof}

\subsection{Nilpotence and compactness}

\begin{proof}[Proof of Theorem~\ref{thm:peripheral-lifting}]
First observe that a continuous map from $\overline\Lambda$ to a
finite group has nilpotent image of $U$. Indeed its kernel on $U$
is an open subgroup containing $R=\bigcap_n\gamma_n(U)$.
Compactness of $U$ and the descending closed family $\gamma_n(U)$
imply that this kernel contains some $\gamma_n(U)$. Conversely,
any finite source homomorphism whose image of $U$ is nilpotent kills
$R$ and factors through $\overline\Lambda$.

Let $q:\Lambda\twoheadrightarrow E$ be a finite quotient. In the
finite group $q(U)$ the lower central series stabilizes. Put
\[
 A=\gamma_\infty(q(U))
       =\bigcap_{n\geq1}\gamma_n(q(U)).
\]
The subgroup $A$ is normal in $E$, and $q(U)/A$ is nilpotent.
Apply Lemma~\ref{lem:finite-peripheral-refinement} to obtain
$q^\dagger:\Lambda\twoheadrightarrow D^\dagger$.
The subgroup $q^\dagger(U)$ is a central extension of $q(U)/A$,
so it is nilpotent as well: a central extension of a group of class
$c$ has class at most $c+1$. It follows that $q^\dagger$ factors
as $\overline q^\dagger\pi$. Define
\[
 q_1^\dagger=\overline q^\dagger\alpha\pi,
 \qquad r_i=\overline q^\dagger(\overline c_i).
\]
This is an epimorphism with the required three conjugacies. The
finite lemma supplies an epimorphism $q_1:\Lambda\to E$ and
conjugators $C_i\in E$ whose images modulo $A$ are the prescribed
images of $\overline c_i$.

We use these finite lifts simultaneously. In the compact space
$\Lambda^3$, impose the following closed conditions on $(c_1,c_2,c_3)$:
\begin{enumerate}
 \item the elements $c_ix_ic_i^{-1}$ have product one;
 \item their first two images generate every finite quotient of $\Lambda$;
 \item the image of each $c_i$ in every finite quotient of
       $\overline\Lambda$ agrees with the image of $\overline c_i$.
\end{enumerate}
For the first condition it is enough to impose its images in finite
quotients. Each finite list of conditions is controlled by one finite
quotient $q:\Lambda\to E$: take the joint image of all source
quotients in the list, including the quotients that factor through
$\overline\Lambda$.

For this joint quotient, $A=\gamma_\infty(q(U))$ maps trivially to
each quotient in the finite list of $\overline\Lambda$, by the nilpotence
observation above. Thus the conjugators $C_i\in E$ from the finite
lemma have all the prescribed images. Arbitrary lifts of the
$C_i$ to $\Lambda$ satisfy the finite list: their peripheral images
have product one in $E$, generate $E$, and have the prescribed
conjugator values in each quotient in the finite list. The family of closed
conditions therefore has the finite intersection property.
Compactness gives a triple satisfying every condition.

The source presentation now defines a continuous endomorphism
$\Phi$ taking $x_i$ to $c_ix_ic_i^{-1}$. In the modular case the
square and cube relations hold because these are conjugates of the
original generators; the product condition gives the required third
image in both cases. Its image maps onto every finite quotient of
$\Lambda$ by the generation conditions. The image is compact and
therefore closed, so $\Phi$ is surjective.

Finally a finitely generated profinite group is Hopfian. To recall the
argument, for each positive integer $d$ there are only finitely many
open normal subgroups of index $d$. Inverse image under a surjective
endomorphism injects this finite set into itself, preserving the index,
and is therefore a permutation. Every such subgroup is consequently
the inverse image of an open normal subgroup and contains the kernel
of the endomorphism. Intersecting all open normal subgroups shows that
the kernel is trivial. Hence $\Phi$ is an automorphism.

The prescribed conjugator conditions imply
$\pi\Phi(x_i)=\alpha\pi(x_i)$ for the two generators, and therefore
$\pi\Phi=\alpha\pi$. If a selected $\overline c_i=1$, then
$c_i\in R$. Replace $\Phi$ by its composite with conjugation by
$c_i^{-1}$ and replace every conjugator $c_j$ by $c_i^{-1}c_j$.
This leaves their images in the quotient and the quotient automorphism
unchanged and makes the selected conjugator equal to one.
\end{proof}
}{}
\section{Why one finite nilpotent quotient is insufficient}
\label{sec:peripheral-finite-obstruction}

We exhibit a finite group generated by a product-one triple and a normal
subgroup whose nilpotent residual does not suffice to correct prescribed
conjugator cosets. The quotient automorphism in the example is the identity,
and the conjugated first two elements still generate the whole group.
The obstruction is an explicit central sign in a central product of
$\operatorname{SL}_2(\F_5)$ and $\operatorname{SL}_2(\F_3)$. A quaternion quotient detects the same
obstruction, explaining why the example cannot arise from the corresponding
full pronilpotent quotient of a free profinite group.

For a group $H$, write $\gamma_{\infty}(H)=\bigcap_{n\geq1}\gamma_n(H)$, with
closed lower central terms when $H$ is profinite. If $D\leq G$, write
$g^D=\{dgd^{-1}:d\in D\}$.

\begin{theorem}\label{thm:peripheral-finite-obstruction}
There exist a finite group $G=\langle g_1,g_2\rangle$, a normal subgroup
$N\trianglelefteq G$, and elements $a_1,a_2,a_3\in G$ with the following
properties. Put $g_3=(g_1g_2)^{-1}$, $D=\gamma_{\infty}(N)$, and
$h_i=a_ig_ia_i^{-1}$. Then
\[
|G|=1440,\qquad |N|=480,\qquad |D|=120,\qquad G/D\cong A_4,
\]
and
\[
h_iD=g_iD\quad(1\leq i\leq3),\qquad
G=\langle h_1,h_2\rangle,\qquad
1\notin h_1^D h_2^D h_3^D.
\]
One may take $a_2=a_3=1$. In particular, the finite quotient datum is an
identity automorphism with specified peripheral conjugators, but those
conjugator cosets cannot be corrected by $D$.
\end{theorem}

\begin{lemma}\label{lem:peripheral-central-sign}
In $S=\operatorname{SL}_2(\F_5)$, put
\[
A=\begin{pmatrix}0&1\\4&1\end{pmatrix},\qquad
B=\begin{pmatrix}0&1\\4&3\end{pmatrix},\qquad
C=\begin{pmatrix}2&2\\1&4\end{pmatrix}.
\]
Then $S=\langle A,B\rangle$, $ABC=I$, and
\[
I\in A^S B^S C^S,\qquad -I\notin A^S B^S C^S.
\]
\end{lemma}

\begin{proof}
We first establish generation, then exclude the opposite central sign.
Direct multiplication gives $C=(AB)^{-1}$ and
\[
(A^{-1}B)^2=\begin{pmatrix}1&1\\0&1\end{pmatrix},\qquad
(AB^{-1})^2=\begin{pmatrix}1&0\\1&1\end{pmatrix}.
\]
The elementary matrices on the right generate $\operatorname{SL}_2(\F_5)$ by row
operations, proving generation and the first inclusion.

Suppose $uvw=-I$, with $u\in A^S$, $v\in B^S$, and $w\in C^S$.
Simultaneous conjugation permits $u=A$. Since
$\operatorname{tr} A=\operatorname{tr} C=1$ and $\operatorname{tr} B=-2$, the equation forces
$\operatorname{tr}(Av)=\operatorname{tr}(-w^{-1})=-1$. Write
\[
v=-I+\begin{pmatrix}r&s\\t&-r\end{pmatrix}.
\]
The determinant condition gives $r^2+st=0$, while the trace condition
gives $t-s-r=0$. Hence $r^2+rs+s^2=0$. This quadratic form has no
nonzero zero over $\F_5$, because its discriminant $-3=2$ is a
nonsquare. Thus $r=s=t=0$ and $v=-I$, contradicting $v\in B^S$.
\end{proof}

\begin{proof}[Proof of Theorem~\ref{thm:peripheral-finite-obstruction}]
We place the sign obstruction in a central product whose nilpotent
residual is the first factor. Let $S,A,B,C$ be as in Lemma~\ref{lem:peripheral-central-sign}. The group $S$ is perfect:
conjugation by $\operatorname{diag}(2,3)$ multiplies the parameter in
each elementary unipotent subgroup by $4$, so the elementary matrix
of parameter $1$ is the commutator with the elementary matrix of
parameter $2$. Since these elementary matrices generate $S$, we have
$[S,S]=S$.

Over $\F_3$, put
\[
E=\operatorname{SL}_2(\F_3),\qquad
i=\begin{pmatrix}0&1\\2&0\end{pmatrix},\qquad
t=\begin{pmatrix}1&1\\0&1\end{pmatrix},\qquad
j=\begin{pmatrix}1&1\\1&2\end{pmatrix}.
\]
The matrices satisfy $i^2=j^2=-I$ and $ji=-ij$. The subgroup
\[
Q=\{\pm I,\pm i,\pm j,\pm ij\}
\]
is quaternion of order eight and is preserved by conjugation by $t$.
Moreover $i,t$ generate $E$: $t$ is an elementary upper unipotent, and
conjugation by $i$ supplies a lower unipotent. Thus $Q$ is normal in
$E$, and $E/Q$ has order three. In particular, $E$ is solvable. The
quotients $E/\{\pm I\}$ and $Q/\{\pm I\}$ are $A_4$ and its normal
Klein four subgroup.

Form
\[
G=(S\times E)/\langle(-I,-I)\rangle.
\]
Use brackets for its elements. Let $D$ be the image of $S\times\{I\}$
and $N$ the image of $S\times Q$. The embedding $S\to D$ is
injective. Since $|S|=120$ and $|E|=24$, the displayed group orders
follow. Also $G/D\cong A_4$ and $N/D\cong V_4$. The identities
$[S,S]=S$ and $[Q,Q]=\{\pm I\}$ give
\[
[N,N]=D,\qquad [D,N]=D,
\]
so $\gamma_{\infty}(N)=D$.

Define
\[
g_1=[(A,i)],\qquad g_2=[(B,t)],\qquad
g_3=[(C,(it)^{-1})].
\]
Their product is one. The subgroup of $S\times E$ generated by
$(A,i),(B,t)$ projects onto each factor. Any common quotient of the
two factors is both perfect and solvable, hence trivial. The elementary
subdirect-product argument therefore makes this subgroup all of
$S\times E$, proving $G=\langle g_1,g_2\rangle$.

Take $a_1=[(I,j)]$ and $a_2=a_3=1$. Then
\[
h_1=[(A,-i)]=[(-A,i)],\qquad h_2=g_2,\qquad h_3=g_3.
\]
The images of $i,j$ commute in $Q/\{\pm I\}$, so $h_iD=g_iD$ for
each $i$. Moreover $B^5=-I$, so $-A,B$ generate $S$; the same
subdirect-product argument proves $G=\langle h_1,h_2\rangle$.

Every element $d_k\in D$ has a unique representative $(s_k,I)$.
The product of the three resulting $D$-conjugates of the $h_k$ is
\[
[( -uvw,I)],\qquad u\in A^S,\quad v\in B^S,\quad w\in C^S.
\]
This is the identity in $G$ precisely when $uvw=-I$: its second
coordinate is $I$, so it cannot represent the nonidentity element
$(-I,-I)$ of the central subgroup being factored out. Lemma~\ref{lem:peripheral-central-sign}
excludes this possibility.
\end{proof}

The full pronilpotent quotient retains more information than the quotient
$G/D$ in Theorem~\ref{thm:peripheral-finite-obstruction}. The following observation makes the
distinction explicit.

\begin{proposition}\label{prop:peripheral-full-quotient-excludes}
Let $\Lambda=\widehat F_2=\langle x_1,x_2\rangle$, put
$x_3=(x_1x_2)^{-1}$, and let $q:\Lambda\twoheadrightarrow G$ send
$x_1,x_2$ to $g_1,g_2$. Set $U=q^{-1}(N)$, $R=\gamma_{\infty}(U)$, and
$\bar\Lambda=\Lambda/R$, and write $\bar x_i$ for the image of $x_i$.
There is no homomorphism
$\alpha:\bar\Lambda\to\bar\Lambda$ and conjugators $\bar c_i$ such
that
\[
\alpha(\bar x_i)=\bar c_i\bar x_i\bar c_i^{-1}
\]
and the images of the $\bar c_i$ in $G/D$ are $a_1D,1,1$.
\end{proposition}

\begin{proof}
Continuous maps of profinite groups commute with the nilpotent residual
when the target is finite: images commute with each lower central term,
and compactness commutes the image with their descending intersection.
Hence $q(R)=\gamma_{\infty}(N)=D$.

There is also a homomorphism
\[
e:\Lambda\longrightarrow E,\qquad e(x_1)=i,\quad e(x_2)=t.
\]
Its composite to $E/\{\pm I\}$ agrees with the composite of $q$ to
$G/D$. Therefore $e(U)=Q$, which is nilpotent, and $e(R)=1$. Write
$\bar e:\bar\Lambda\to E$ for the resulting homomorphism.

The asserted finite conjugator images would imply
\[
\bar e(\bar c_1)\in\{j,-j\},\qquad
\bar e(\bar c_2),\bar e(\bar c_3)\in\{I,-I\}.
\]
Consequently the images under $\bar e\alpha$ of the three peripheral
generators would be $-i$, $t$, and $(it)^{-1}$. Their product is $-I$,
contradicting $x_1x_2x_3=1$.
\end{proof}

\IfFileExists{characters.tex}{\section{From even weights to full profinite characters}\label{sec:characters}
This section explains why rational even-weight indecomposables produce
characters of the full gentle group which vanish on the pentagon group.
The rational argument does not identify a mod-$p$ pentagon kernel with a
reduction of its characteristic-zero kernel.

We use the Deligne--Drinfeld theorem in the precise form proved by
OpenAI and streamlined by Willwacher \cite[Theorem 1.1 and the discussion
of (7), pp.~1--2]{Willwacher}: $\grt(\Q)$ is generated, as a graded Lie
algebra, by one element in each odd weight at least three.  Brown's
injectivity theorem supplies freeness, although generation alone is enough
here.  As $\grt\subset W$, it follows that
\begin{equation}
 \ker P_m=\grt(\Q)_m\subset[W,W]_m
       \qquad(m\text{ even}).\label{eq:even-pentagon}
\end{equation}
Indeed an even-weight Lie word in odd-weight generators has at least two
letters and is a bracket of two positive-weight Lie words.  The same
inclusion holds after scalar extension to $\Q_p$.

\begin{theorem}[Even-weight characters]\label{thm:even-characters}
Let $p\geq5$ be prime, let $m\geq2$ be even, and put
$h_m=\dim_\Q(W/[W,W])_m$.  For every $k\geq1$ such that
$(p-1)p^{k-1}$ divides $m$, there is a continuous surjection
\[
 \GTz/\overline{\langle\!\langle\GT\rangle\!\rangle}
       \longrightarrow(\Z/p^k\Z)^{h_m}.
\]
The map is already onto on the global multiplier-one subgroup of $\GTz$.
In particular weight sixteen gives a quotient $(\Z/17\Z)^{38}$.
\end{theorem}

The proof needs a full realization theorem.  We spell out its local
part, and use Theorem~\ref{thm:peripheral-lifting} for the passage to
the full profinite group.

\begin{lemma}[Realization of special directions]\label{lem:realization}
For $p\geq5$, let $\mathcal U$ be the multiplier-one subgroup in the
image of $\GTz$ on the free pro-$p$ group $P$ on $X,Y$.
Let $\mathcal U_n$ consist of pairs whose normalized conjugator belongs
to $\gamma_nP$.  Then
\[
 \mathcal U_2=\mathcal U_3=\mathcal U,
 \qquad \mathcal U_n/\mathcal U_{n+1}\simeq W_n(\Z_p),
 \qquad [\mathcal U_i,\mathcal U_j]\subset\mathcal U_{i+j}.
\]
Every element of the local multiplier-one group has a full gentle lift
of global multiplier one.
\end{lemma}
\begin{proof}
Write $Z=(XY)^{-1}$.  The ordinary lower-central graded Lie algebra of
$P$ is the free Lie algebra over $\Z_p$: integral Hall collection gives
the coordinates in each free nilpotent quotient, and every finite
$p$-group is nilpotent.  A leading normalized conjugator gives
$D_q(x)=0$, $D_q(y)=[y,q]$.  Preserving the three peripheral classes
and commuting with their permutations gives precisely the equations
defining $W_n$, as proved in Section~\ref{sec:dimensions}.  This also
proves injectivity of the leading map.

Conversely, start with an invariant outer special derivation of degree
$n$.  Normalize it to fix $x$, and lift its conjugators at $y,z$ to
$\gamma_nP$.  The prospective images of $X,Y,Z$ have product error in
$\gamma_{n+2}P$.  If the error lies in $\gamma_jP$, changing the two
conjugators in degree $j-1$ changes its degree-$j$ term by
\[
 (a,b)\longmapsto[a,y]+[b,-x-y].
\]
This map is onto: its image is $[L_{j-1},x]+[L_{j-1},y]=L_j$.
Successively cancel every error.  The corrections converge because
every finite $p$-group quotient kills sufficiently deep lower-central
terms.  The resulting endomorphism is the identity on the Frattini
quotient, hence is onto by the pro-$p$ basis theorem and is an
automorphism by finite generation.

The special outer automorphisms with trivial Frattini action form a
pro-$p$ group.  Their filtration is stable under the geometric $S_3$.
Fixed points lift through an onto equivariant map of pro-$p$ groups
when the acting finite group has order prime to $p$.  To see this,
the discrepancy of a lift is a nonabelian one-cocycle in the kernel;
the corresponding two complements in its semidirect product are
conjugate.  For a finite $p$-kernel, induction through a central
elementary abelian subgroup and averaging prove this assertion.
For a pro-$p$ kernel, compactness of the closed sets of possible
conjugating elements proves it from the finite cases.  Apply this to
the degree-$n$ quotient.  It changes only higher terms and gives an
$S_3$-invariant special outer automorphism with the required leading term.

Here is the explicit passage to the normalized group equations.
Let $\mathcal B=\widehat{C_2*C_3}$, with generators $s,r$, and set
$\tau=sr$, $F=\ker(\mathcal B\to S_3)$,
$X=\tau^{-2}$, $Y=sXs$.  Schreier rewriting gives $F=\widehat F_2$
on $X,Y$, and conjugation by $r$ cycles $X,Y,Z$.  The quotient
$B_p=\mathcal B/O^p(F)$ has kernel $P$ and outer quotient $S_3$.
The centre of $P$ is trivial, so
\[
 B_p\simeq\operatorname{Aut}(P)
           \mathbin{\times}_{\operatorname{Out}(P)} S_3.
\]
An invariant outer class therefore gives an automorphism of $B_p$
acting trivially on $S_3$.  Choose its representative to fix $X$.
Its image of $\tau$ has the form $X^a\tau$: both elements centralize
$X$, their quotient lies in $P$, and $C_P(X)=\langle X\rangle$.
Equality of their squares implies $2a=0$ in $\Z_p$, so it fixes $\tau$.
If $f_1=s\phi(s)$, then $f_1$ has abelianization $(a,-a)$.
Conjugating $\phi$ by $X^a$ replaces it by
$f=Y^af_1X^{-a}\in P'$.  The resulting formulas are
\[
 \phi(s)=sf,\qquad \phi(r)=f^{-1}r,
 \qquad X\mapsto X,\quad Y\mapsto f^{-1}Yf.
\]
The relations $s^2=r^3=1$ give, in $P$,
\[
 f(Y,X)f(X,Y)=1,\qquad f(Z,X)f(Y,Z)f(X,Y)=1.
\]
These are the multiplier-one gentle equations.  Centralizers in
$P/\gamma_{n+1}P$ show that the normalization changes no degree-$n$
outer term: an inner change fixing $x$ has zero homogeneous
conjugator in degree $n>1$.

Use this automorphism in the $p$-component of the maximal pronilpotent
pure quotient of $\mathcal B$, and the identity in every other component.
It preserves the positive $s$- and $r$-classes: their complements over
$C_2,C_3$ are conjugate by the same coprime argument.  It fixes $\tau$
exactly.  The full peripheral lifting theorem gives an automorphism
of $\mathcal B$ with these prescribed conjugators.  An inner correction
from the kernel makes it fix $\tau$ exactly.  Its normalized $f=s\phi(s)$
has zero abelianization at every prime, and the same torsion relations
give the full gentle equations.  Thus it has global multiplier one and
the prescribed local image.  This argument applies to an arbitrary
local solution, not only to a leading term.

Finally, composition of normalized conjugators gives the Ihara bracket
on the associated graded, hence the displayed commutator inclusion.
Normalization removes degree one and $W_2=0$ removes degree two.
\end{proof}

\begin{lemma}[Rational Lie algebra of a compact unipotent group]
\label{lem:compact-log}
Let $U$ be a closed subgroup of a unipotent matrix group over $\Q_p$.
Set $L=\operatorname{span}_{\Q_p}\log U$.  Then $L$ is a nilpotent Lie
algebra and $U$ contains $\exp(J)$ for some full $\Z_p$-lattice
$J\subset L$.  These statements apply to every closed subgroup of $U$.
\end{lemma}
\begin{proof}
All logarithms, exponentials and Baker--Campbell--Hausdorff expressions
are finite polynomials.  For $X=\log u$ and $Y=\log v$, the polynomial
$\operatorname{BCH}(aX,bY)$ belongs to $L$ for nonnegative integers
$a,b$.  Evaluating on a sufficiently large integer grid and inverting
the rational Vandermonde matrices shows that every coefficient belongs
to $L$.  Its $ab$-coefficient is $[X,Y]/2$, proving closure under brackets.

Choose a basis $X_i=\log u_i$ of $L$ from logarithms of elements of $U$.
For $t_i\in\Z_p$, continuity gives $u_i^{t_i}=\exp(t_iX_i)\in U$.
In these basis coordinates,
\[
 (t_i)\longmapsto\log(u_1^{t_1}\cdots u_d^{t_d})=t+R(t),
\]
where $R$ has no constant or linear term.  On a sufficiently small ball
$p^r\Z_p^d$, $R$ has Lipschitz constant at most $p^{-1}$ and maps that
ball into $p^{r+1}\Z_p^d$.  For each $y$ in the ball, contraction of
$t\mapsto y-R(t)$ solves $t+R(t)=y$.  Thus the exponential of the
corresponding full lattice is contained in $U$.  The proof uses only
closedness and applies verbatim to a closed subgroup.
\end{proof}

The ordinary lower-central Lie ring of the pure braid group on four
strands is the integral infinitesimal braid Lie ring
\cite[\S4, p.~223]{CohenFalkRandell}.  This remains true after
pro-$p$ completion and extension to $\Z_p$: Artin's successive free-group
decomposition has actions trivial on abelianizations, and Hall collection
replaces the integral coordinates in each nilpotent quotient by
$p$-adic coordinates.  Consequently, if the normalized conjugator of
a pentagon element first appears in weight $n$, the leading term of
its group pentagon is exactly $P_n(q)=0$ in
$\mathfrak t_{4,n}(\Z_p)$.  This is an inclusion of group leading terms
in the integral kernel; no equality with a reduced rational kernel is
needed.

\begin{proof}[Proof of Theorem~\ref{thm:even-characters}]
Let $\mathcal G$ be the local pro-$p$ image of $\GTz$, and set
\[
 G=\mathcal G/\mathcal U_{m+1},\qquad
 U=\mathcal U/\mathcal U_{m+1}.
\]
The action on the integral Magnus algebra through augmentation degree
$m+1$ realizes $U$ faithfully by unipotent matrices.  Put
$L=\operatorname{span}_{\Q_p}\log U$.  Lemmas
\ref{lem:realization} and \ref{lem:compact-log} identify its degree
filtration with
\[
 L_n/L_{n+1}=W_n(\Q_p),\qquad 3\leq n\leq m.
\]
Indeed every leading group term gives the indicated Lie term, and the
open exponential lattice ensures that every filtered Lie vector, after
multiplication by a power of $p$, is a logarithm of a filtered group element.

Choose an element $g_0$ in the local Galois image whose multiplier is
$1+p$.  The action $A=\operatorname{Ad}(g_0)$ on $L_n/L_{n+1}$ is
$(1+p)^n$: conjugation of a degree-$n$ derivation by the scalar
action on degree one scales it by that power.  These eigenvalues are
distinct over $\Q_p$.  Applying the factors successively to the
filtration gives
\[
 \prod_{n=3}^{m}(A-(1+p)^n)=0.
\]
Thus $A$ is semisimple, and its eigenspaces give a graded splitting
\[
 L=\bigoplus_{n=3}^{m}L^{(n)},\qquad L^{(n)}\simeq W_n(\Q_p).
\]
Because $A$ is a Lie automorphism, brackets respect this grading and
are the Ihara brackets on the leading terms.

Choose the full space of degree-$m$ linear functionals on
$L/[L,L]$, extending each by zero on the other eigenspaces.  They give
a map $\ell:L\to\Q_p^{h_m}$ annihilating $[L,L]$.
Consequently $u\mapsto\ell(\log u)$ is an additive homomorphism.
Its image is a compact $\Z_p$-submodule spanning $\Q_p^{h_m}$, and
hence is a full free lattice.  Choosing a basis of this image gives a
continuous surjection
\[
 \Phi:U\longrightarrow\Z_p^{h_m}.
\]
Conjugation on $L/[L,L]$ factors through
$G/U=\Z_p^\times$, so every such conjugation commutes with $A$ on
that abelianization.  It preserves its weight spaces.  On the top
weight it acts as on the leading quotient, by the scalar
$\lambda_p(g)^m$.  Therefore
\begin{equation}
 \Phi(gug^{-1})=\lambda_p(g)^m\Phi(u).
 \label{eq:equivariance}
\end{equation}

Let $P_0$ be the image of $\GT$ in $G$ and put $P_1=P_0\cap U$.
By Lemma~\ref{lem:compact-log},
$L_P=\operatorname{span}_{\Q_p}\log P_1$ is a Lie subalgebra and
$P_1$ contains an exponential neighborhood in $L_P$.
Since $g_0$ is Galois, it belongs to $P_0$; thus $L_P$ is
$A$-stable.  The polynomial projector $e_m$ onto $L^{(m)}$ preserves
$L_P$.  For $h\in P_1$, put $X=e_m\log h$.  Then
$X\in L_P\cap L_m$, and for some $r$, $\exp(p^rX)\in P_1$.
This is the image of a full pentagon element whose local conjugator
has no terms below weight $m$.  Its leading term therefore satisfies
the linear pentagon.  By \eqref{eq:even-pentagon}, it is a bracket in
$W_m(\Q_p)$, so $\ell(p^rX)=0$ and $\ell(X)=0$.
This proves $\Phi(P_1)=0$.

Reduce $\Phi$ modulo $p^k$.  The divisibility hypothesis gives
$a^m=1$ in $(\Z/p^k\Z)^\times$ for every $a\in\Z_p^\times$,
so \eqref{eq:equivariance} makes this reduction, denoted $\theta$,
invariant under all of $G$.  The multiplier on $P_0$ is onto,
because it contains the Galois image.  Hence $G=UP_0$.
Define
\[
 \chi(up)=\theta(u)\qquad(u\in U,\ p\in P_0).
\]
Two such factorizations differ by an element of $P_1$, so the value
is well defined.  Invariance of $\theta$ proves the homomorphism law.
The multiplication map from the compact space $U\times P_0$ onto
$G$ is a quotient map, which proves continuity.  The homomorphism
vanishes on $P_0$ and is onto.  Composing with the full gentle action
proves the asserted quotient.  The final clause of
Lemma~\ref{lem:realization} supplies all values on global multiplier one.
\end{proof}

\begin{corollary}[All finite odd-primary orders]\label{cor:all-orders}
Put $Q=\GTz/\overline{\langle\!\langle\GT\rangle\!\rangle}$.
There is a continuous surjection
\[
 Q^{\mathrm{ab}}\longrightarrow
       \prod_{\substack{p\geq5\\p\ \mathrm{prime}}}
          \prod_{k,j\geq1}\Z/p^k\Z.
\]
For every prime $p\geq5$, the vector space
$\Hom_{\cts}(Q,\F_p)$ has dimension exactly $\aleph_0$.
\end{corollary}
\begin{proof}
We first prove that $h_m$ tends to infinity along every even arithmetic
progression.  Let $a_n$ and $b_n$ be as in
\eqref{eq:dimension}--\eqref{eq:raw}, and put $r_n=a_n2^{-n}$.
The identity contribution to the Witt formula gives
\[
 a_n=\frac{2^n}{3n(n+1)}+O(2^{n/2}),\qquad
 r_n\sim\frac1{3n^2}.
\]
All nonidentity contributions and all nontrivial M\"obius divisors
have exponential size at most the stated error.
Write $l_n=\dim L_n$.  The dimension of
$\operatorname{OutSDer}_n$ is $2l_n-l_{n+1}$, whence
\[
 \sum_{n=2}^N r_n\leq
 \sum_{n=2}^N\frac{2l_n-l_{n+1}}{2^n}
       =\frac12-\frac{l_{N+1}}{2^N}.
\]
At $n=5$, the upper-bound dimension is three whereas $a_5=1$.
Consequently $S=\sum_{n\geq2}r_n\leq7/16$.
For any fixed $M$, the terms with $i\leq M$ or $n-i\leq M$ in
$\sum_i r_ir_{n-i}/r_n$ tend to $2\sum_{i\leq M}r_i$.
In the remaining half-interval, $r_{n-i}/r_n$ is bounded uniformly
and $\sum_{i>M}r_i=O(M^{-1})$.  Letting $M$ increase proves
\[
 \frac{b_n}{a_n}\longrightarrow S<1,
 \qquad h_n\geq a_n-b_n\sim(1-S)\frac{2^n}{3n^2}.
\]
The diagonal subtraction in the unordered pair count is exponentially
smaller and does not affect the limit.

Fix $p,k$.  Theorem~\ref{thm:even-characters}, with arbitrarily large
multiples of $(p-1)p^{k-1}$, gives arbitrarily many characters to
$\Z/p^k\Z$ whose reductions modulo $p$ are independent already on
global multiplier one.  Thus the space of these reductions is
infinite-dimensional.  Enumerate all pairs $(k,j)$.  At each stage
choose a character whose modulo-$p$ restriction lies outside the
finite span of the preceding restrictions.  For any finite collection,
let $H$ be its image in the target product $T$ of cyclic $p$-groups.
Independence gives $H+pT=T$.  Hence the finite group $T/H$ equals
$p(T/H)$ and is zero.  The finite product map is onto.

For different primes the primary targets have coprime orders, so a
subdirect product is their product.  Every finite projection of the
displayed infinite map is therefore onto.  Compactness proves
surjectivity onto the whole product.  Finally, the gentle group is
countably based: a finitely generated free profinite group has a
countable cofinal family of characteristic open subgroups.  Its
automorphism group, and the closed gentle subgroup, inherit this
property.  There are only countably many clopen subsets of a compact
space with a countable clopen basis, hence only countably many maps
to $\F_p$.  The lower bound already proved gives equality.
\end{proof}

\begin{remark}
The corollary specifies a quotient of the abelianization.  It does not
identify the whole abelianization, prove that the original quotient is
abelian, or construct a compatible $\Z_p$ quotient.  All its characters
vanish on the full pentagon group, so they give no separation of that
group from the absolute Galois group.
\end{remark}
}{}
\IfFileExists{weight17.tex}{\section{The pentagon in weight sixteen modulo seventeen}
\label{sec:mod17}
The preceding character theorem uses rational compact images.  We now
compute the entire finite-field pentagon quotient in weight sixteen,
and give an integral construction that realizes every functional in that
quotient.  No equality between rational and modular pentagon kernels is
assumed.

\begin{theorem}\label{thm:pentagon17}
Over $\F_{17}$ the pentagon map in weight sixteen has the following ranks:
\[
 \operatorname{rank}P_{16}=70,\qquad
 \dim\ker P_{16}=5,\qquad
 \operatorname{rank}(P_{16}|_{[W,W]_{16}})=32.
\]
Its kernel is contained in $[W,W]_{16}$.  Consequently the induced map
\[
 W_{16}/[W,W]_{16}\longrightarrow
 \mathfrak t_{4,16}/P_{16}([W,W]_{16})
\]
is injective and has rank thirty-eight.  Every functional on
$(W/[W,W])_{16}$ annihilates the linear pentagon kernel.
\end{theorem}
\begin{proof}
We first give the upper bound on the rank.  The infinitesimal braid Lie
algebra is torsion free over $\Z$: its presentation gives the semidirect
decomposition
\[
 \mathfrak t_4=\Lie_\Z(t_{14},t_{24},t_{34})\rtimes\mathfrak t_3,
 \qquad
 \mathfrak t_3=\Z(t_{12}+t_{13}+t_{23})\oplus
                   \Lie_\Z(t_{12},t_{23}).
\]
Indeed the cross-brackets specified by the infinitesimal braid relations
define derivations of the free first factor; checking the relations on
its three generators gives inverse homomorphisms to and from this
semidirect presentation.  Each homogeneous component is therefore free
of finite rank.  Section~\ref{sec:dimensions} proves that
$W_{16}(\Z_{17})$ is free and reduces onto $W_{16}(\F_{17})$.
The rank of its integral pentagon matrix cannot increase modulo seventeen.

Brown's injection gives a free Lie subalgebra of the rational pentagon
kernel with one generator in each odd weight at least three
\cite[Theorem~1.1 and \S1.1]{Brown}.  Its weight-sixteen part has
dimension five.  The free Lie algebra on the odd
generators has three two-letter brackets, of weights $(3,13)$, $(5,11)$
and $(7,9)$, and two four-letter words, of multidegrees $(3,3,3,7)$
and $(3,3,5,5)$.  Each of those latter multidegrees has dimension one;
this follows directly by Jacobi, or by the multigraded Witt formula.
No word with six letters has weight sixteen.  Thus the rational kernel
has dimension at least five, the rational rank is at most $75-5=70$,
and the modular rank is at most seventy.  This upper bound uses Brown's
injection alone; it does not require generation of the whole rational
pentagon Lie algebra.

Here is a complete finite specification of a lower-bound certificate.
All row reductions in this paragraph are over $\F_{17}$, using the
first possible pivot in the stated order.  In the lexicographic Lyndon
basis of $L_{16}$ form the columns of
\[
 \Pi=(1-s)(2-r-r^2),\qquad
 s(q)=q(y,x),\quad r(q)=q(y,-x-y).
\]
The image of $\Pi$ is exactly
$\{q:sq=-q,\ (1+r+r^2)q=0\}$: on the standard representation,
$2-r-r^2$ is three times the identity, while it is zero on the other
two irreducible representations.  Retain the first independent columns
of $\Pi$, obtaining $1360$ polynomials.  Apply the special equation
\eqref{eq:special} to these columns, taking length-seventeen Lyndon-word
coefficients.  Its matrix has rank $1285$.  For each free column, in
increasing order, take the reduced-row-echelon null vector with that free
coordinate one and the other free coordinates zero.  This gives the
ordered basis $q_1,\ldots,q_{75}$ of $W_{16}$ used below.  The equations
on all rows vanish on these vectors; independence and the dimension
formula prove that no solution is omitted.

For every output word in Table~\ref{tab:rows17}, take the indicated
coefficient of $P_{16}(q)(a)$ in the three-generator representation of
Section~\ref{sec:cert}.  List the left column of words first and the
right column second, and call the resulting $70\times75$ matrix $R$.
The coefficient recurrence \eqref{eq:contraction}, split into blocks of
eight operators, gives
\begin{equation}
 \det R_{\{1,\ldots,70\},\{1,\ldots,70\}}=4\quad\text{in }\F_{17}.
 \label{eq:minor17}
\end{equation}
This proves rank at least seventy and hence equality.  In particular,
these coefficient rows have exactly the same kernel as $P_{16}$.

Let $b_1,\ldots,b_{37}$ be the ordered bracket basis of
Section~\ref{sec:brackets}, obtained by omitting the twenty-eighth of the
thirty-eight columns.  The $37\times5$ matrix $X$ in
Table~\ref{tab:kernel17} defines five vectors
$v_j=\sum_i X_{ij}b_i$.  Applying the same coefficient recurrence gives
\begin{equation}
 R(v_1\ \cdots\ v_5)=0.
 \label{eq:kernel17}
\end{equation}
The minor of $X$ on rows $1,2,3,22,30$ has determinant $5$ modulo
seventeen, so the five vectors are independent.  They therefore span
$\ker R=\ker P_{16}$ and all lie in the bracket space.  Its dimension
is thirty-seven, so the restriction of $P_{16}$ to it has rank
$37-5=32$.  The quotient assertion now follows from elementary linear
algebra.  Equations \eqref{eq:minor17}--\eqref{eq:kernel17}, together
with the deterministic basis prescription, are finite certificates of
every modular assertion used in the proof.
\end{proof}

\begin{table}[p]
\centering\small
\begin{tabular}{rllrll}\toprule
row&input&output word&row&input&output word\\\midrule
1&$a$&\texttt{abcbbaaaaaacaaaca}&36&$a$&\texttt{baacabaaaacacabba}\\
2&$a$&\texttt{ccabbbcacacbbabaa}&37&$a$&\texttt{babbaaaacbccaaaaa}\\
3&$a$&\texttt{aabaccabbbcaccabb}&38&$a$&\texttt{cabaabccabbaccbaa}\\
4&$a$&\texttt{baabcccaaaabbbaca}&39&$a$&\texttt{cbaaaccbcacbbabba}\\
5&$a$&\texttt{accbacbaacaabbacb}&40&$a$&\texttt{abaaaaabcacacabaa}\\
6&$a$&\texttt{cbcabacbabcaaabaa}&41&$a$&\texttt{aacbaaaaacaabbbac}\\
7&$a$&\texttt{cabaabaaabacaaaca}&42&$a$&\texttt{acabaacacabaaaaba}\\
8&$a$&\texttt{abaabbacaaaaaacca}&43&$a$&\texttt{abaabccacababbaca}\\
9&$a$&\texttt{aaaaabbbaaccacbbc}&44&$a$&\texttt{abacaacaabaaaacba}\\
10&$a$&\texttt{abbbacbaccacabaca}&45&$a$&\texttt{cccababcabacbabba}\\
11&$a$&\texttt{aaaaaaaacaabcbcab}&46&$a$&\texttt{aabcaaacaaaabbcaa}\\
12&$a$&\texttt{aabacccacaacbbbba}&47&$a$&\texttt{bcaaabccabbbccbaa}\\
13&$a$&\texttt{caacabcaaabcbbaaa}&48&$a$&\texttt{aabacacabaaabcaba}\\
14&$a$&\texttt{bcaaccbaabccbbbaa}&49&$a$&\texttt{aaacacbbaaacbaaaa}\\
15&$a$&\texttt{acbcaacacbabcaabb}&50&$a$&\texttt{accbaccbabcabaaba}\\
16&$a$&\texttt{bcaaccbcabaaabaaa}&51&$a$&\texttt{bbccaabacabcbbcaa}\\
17&$a$&\texttt{caabbbbcabccabcaa}&52&$a$&\texttt{abaaabcbcacacaaab}\\
18&$a$&\texttt{aaaaccaaaabacbbba}&53&$a$&\texttt{abacbbaacaacaccbb}\\
19&$a$&\texttt{aaaaaacabccabacbb}&54&$a$&\texttt{acbaacbcaabcababa}\\
20&$a$&\texttt{baaaabaaaaccbbcaa}&55&$a$&\texttt{acccbababcacabbba}\\
21&$a$&\texttt{abbaabcaaacaaacab}&56&$a$&\texttt{cacbaababbaaaacaa}\\
22&$a$&\texttt{cbbaacbaacabbcaaa}&57&$a$&\texttt{abaaaabacbbaccaaa}\\
23&$a$&\texttt{acabacbaaaacbaaaa}&58&$a$&\texttt{caabbaaaacbaaabca}\\
24&$a$&\texttt{cbbaabacbccbcabaa}&59&$a$&\texttt{aaaaaaacbacbabcaa}\\
25&$a$&\texttt{abbcbcbbacacaaabc}&60&$a$&\texttt{abaaaacbbcccbbcaa}\\
26&$a$&\texttt{bbacbacbaabbaccca}&61&$a$&\texttt{acaacabbbccbaaaaa}\\
27&$a$&\texttt{acbaccabaaabcaaba}&62&$a$&\texttt{aacbbcabcaaaaaaba}\\
28&$a$&\texttt{aabaaaaacbacabaca}&63&$a$&\texttt{aaaababaaaaabccca}\\
29&$a$&\texttt{aaaacbaabaacacbbc}&64&$a$&\texttt{bcaabbaaaaacaaaca}\\
30&$a$&\texttt{abcbaccaaababacab}&65&$a$&\texttt{aabacbccabbaaccba}\\
31&$a$&\texttt{aaaabbacbaaacaaca}&66&$a$&\texttt{cccababacabcbbaba}\\
32&$a$&\texttt{acbaabbccabacaaba}&67&$a$&\texttt{aaaacbcaabaacbaba}\\
33&$a$&\texttt{ccaaabbcaaabcbaba}&68&$a$&\texttt{aacbbaabaacaaabac}\\
34&$a$&\texttt{acaababcaaaabaaac}&69&$a$&\texttt{abacaabcbabbcacbc}\\
35&$a$&\texttt{ccaaaccaaababbbaa}&70&$a$&\texttt{aabaaaaacbacbacba}\\
\bottomrule\end{tabular}
\caption{The seventy coefficient rows of the weight-sixteen pentagon certificate over $\F_{17}$.}\label{tab:rows17}
\end{table}

\begin{table}[!htbp]
\centering\small
\begin{tabular}{r|rrrrr@{\qquad}r|rrrrr}\toprule
$i$&\multicolumn{5}{c}{$(X_{i1},\ldots,X_{i5})$}&$i$&\multicolumn{5}{c}{$(X_{i1},\ldots,X_{i5})$}\\\midrule
1&7&1&3&10&15&20&8&10&14&16&14\\
2&3&10&12&14&1&21&6&8&11&16&0\\
3&15&11&6&15&6&22&11&4&5&5&8\\
4&16&16&12&5&5&23&7&1&14&14&2\\
5&5&3&4&8&2&24&5&8&10&10&16\\
6&8&11&7&16&1&25&12&9&7&7&1\\
7&9&5&0&1&12&26&9&11&1&1&5\\
8&10&11&9&1&5&27&16&12&15&15&7\\
9&7&0&6&6&15&28&15&7&13&13&14\\
10&7&3&3&7&9&29&16&12&15&15&7\\
11&0&6&11&2&5&30&11&15&7&2&11\\
12&10&12&15&14&13&31&0&0&0&0&0\\
13&9&7&6&4&0&32&12&4&3&13&12\\
14&9&7&11&9&12&33&8&14&2&3&8\\
15&8&1&14&4&7&34&15&5&8&12&15\\
16&14&3&6&10&10&35&0&0&0&0&0\\
17&9&11&11&3&0&36&4&7&1&10&4\\
18&6&4&11&5&12&37&14&16&12&1&14\\
19&2&1&2&9&1&&&&&&\\
\bottomrule\end{tabular}
\caption{Bracket coordinates of five independent pentagon-kernel vectors. All entries are residues modulo seventeen.}\label{tab:kernel17}
\end{table}


\begin{theorem}[Integral top-weight characters]\label{thm:topweight}
For every prime $p\geq5$, put $m=p-1$.  After a choice of tame splitting,
there is an injective linear map
\[
 \bigl(W_m(\F_p)/([W,W]_m(\F_p)+\ker P_m)\bigr)^*
 \longrightarrow
 \Hom_{\cts}\bigl(\GTz/\overline{\langle\!\langle\GT\rangle\!\rangle},
                  \F_p\bigr).
\]
The character associated with $\ell$ restricts to $\ell$ on the leading
weight-$m$ subgroup.  The characters are independent already on the
global multiplier-one subgroup.  In particular the integral certificate
in weight twelve gives a nonzero character of order thirteen without
using the Deligne--Drinfeld theorem; Theorem~\ref{thm:pentagon17} realizes
all thirty-eight weight-sixteen functionals as characters of order
seventeen.
\end{theorem}
\begin{proof}
Use $G,U,P_0,P_1$ as in the proof of
Theorem~\ref{thm:even-characters}, with $m=p-1$.
The group $U$ is a finitely generated torsion-free pro-$p$ group of
nilpotency class at most $\lfloor m/3\rfloor<p$.
The Lazard correspondence therefore gives a finite free
$\Z_p$-Lie algebra $\mathcal L=\log U$, compatible with closed subgroups,
normal subgroups and automorphisms
\cite[Theorem 1(c)--(f)]{LazardTransport}.
Its filtration has associated graded $W_n(\Z_p)$ for $3\leq n\leq m$.
This can also be read directly from the finite logarithm series: its
denominators are prime to $p$, and its first term is unchanged on each
leading filtered quotient.  In particular
$\log P_1$ is a closed $\Z_p$-Lie subalgebra of $\mathcal L$.

Let $\omega\in\Z_p^\times$ generate the group of $(p-1)$st roots of
unity.  The preimage of this finite group in $P_0$ is an extension by
the pro-$p$ group $P_1$.  Choose a lift $g$ of $\omega$.  The sequence
$g^{p^j}$ converges to an element $t\in P_0$ of order $p-1$ and
multiplier $\omega$: in each finite quotient, the order of $g$ divides
$(p-1)p^a$ for some $a$, and these powers are eventually constant, being zero
modulo $p^a$ and one modulo $p-1$.  Its integral eigenspace projectors are
\[
 e_n=\frac1{p-1}\sum_{j=0}^{p-2}\omega^{-nj}\operatorname{Ad}(t)^j,
        \qquad 3\leq n\leq p-1.
\]
These weights have distinct residues modulo $p-1$, and the filtration
shows that
\[
 \mathcal L=\bigoplus_{n=3}^{m}\mathcal L^{(n)},\qquad
 \mathcal L^{(n)}=e_n\mathcal L\simeq W_n(\Z_p).
\]
Brackets respect ordinary weight: if $i+j\leq m$, their leading term
has eigenvalue $\omega^{i+j}$ and there is no other eigenspace with that
residue; if $i+j>m$, the filtration makes the bracket zero.
Thus the top component of $[\mathcal L,\mathcal L]$ is exactly the
integral Ihara bracket span.  In particular its reduction lies in
$[W,W]_m(\F_p)$, with no saturation assertion required.

For a functional $\ell$ in the theorem define
\[
 \theta_\ell(u)=\ell\bigl(e_m\log u\bmod p\bigr).
\]
The Baker--Campbell--Hausdorff corrections are brackets with
$p$-integral coefficients.  They vanish under $\ell$, so $\theta_\ell$
is a continuous homomorphism $U\to\F_p$.
Conjugation on the integral Lie abelianization factors through
$G/U=\Z_p^\times$ and commutes there with $\operatorname{Ad}(t)$.
On its top summand it acts by $\lambda_p^m$; this is one modulo $p$.
Hence $\theta_\ell$ is invariant under all of $G$.

Since $t\in P_0$, every projector $e_n$ preserves $\log P_1$.
For $h\in P_1$, the element
$\exp(e_m\log h)$ therefore belongs to $P_1$ and has no terms below
weight $m$.  Its leading term satisfies the integral linear pentagon.
Its reduction belongs to $\ker P_m$ over $\F_p$, and $\ell$ kills it.
Thus $\theta_\ell(P_1)=0$.  As $G=UP_0$, the well-defined continuous
homomorphism $\chi_\ell(up)=\theta_\ell(u)$ extends it to $G$ and
vanishes on $P_0$, exactly as in the rational proof.

The construction is linear in $\ell$.  On $U_m$ it is precisely the
prescribed functional of the leading coefficient.  Realization and
full peripheral lifting give every such coefficient on global
multiplier one, proving injectivity and the final assertions.
\end{proof}

\begin{corollary}[Strictness inside a quadratic kernel]
\label{cor:quadratic-kernel}
Let $\delta:\GTz\to\F_2$ be a continuous character.  Its kernel
contains a global multiplier-one element outside $\GT$.  If $\delta$
is nonzero and vanishes on $\GT$, then
\[
 \GT\varsubsetneqq\ker\delta\varsubsetneqq\GTz.
\]
This conclusion does not use the Deligne--Drinfeld theorem.
\end{corollary}
\begin{proof}
The weight-twelve certificate and Theorem~\ref{thm:topweight} give a
character $\chi:\GTz\to\F_{13}$ which vanishes on $\GT$ and is onto
on global multiplier one.  Choose such an element $g$ with $\chi(g)=1$.
Then $\delta(g^2)=0$ and $\chi(g^2)=2\ne0$, so $g^2$ is the required
element.  The two strict inclusions follow from the additional hypotheses
on $\delta$.
\end{proof}
}{}
\IfFileExists{quadratic_sign_lifting.tex}{\section{Quadratic characters from a finite cover}
\label{sec:quadratic-beta-characters}

The quadratic construction uses a cover of degree thirty-six of the
three-punctured sphere. Its homology separates the cubic character pair
from the other sixth-root characters. An involution on that pair lifts
through every finite central extension, which is the additional fact
needed to pass from the finite homology calculation to a full peripheral
automorphism.

Use the modular marking
\[
 \Gamma=\widehat{C_2*C_3}=\langle s,r\mid s^2=r^3=1\rangle,
 \qquad \tau=sr,
 \qquad F=\ker(\Gamma\to S_3),
\]
and choose $x=\tau^{-2}$ and $y=sxs$ as free generators of $F$.
Set $z=(xy)^{-1}$. Conjugation by $s$ interchanges $x,y$, and
conjugation by $r$ cyclically permutes $x,y,z$. Put
\[
 U_6=\overline{F'F^6},\qquad
 C=F/U_6\simeq(\mathbb Z/6\mathbb Z)^2,
 \qquad \Gamma_6=\Gamma/U_6.
\]
Here $F'$ denotes the closed commutator subgroup and $F^6$ the closed
subgroup generated by sixth powers. The group $\Gamma_6$ has order
$216$ and abelianization $C_6$: the inclusion $U_6\subset\Gamma'$
follows because $x,y$ have order dividing three in $\Gamma^{\mathrm{ab}}$.

\subsection{The homology of the cover}

For a prime $p\geq5$, write
\[
 M_p=U_6/\overline{U_6'U_6^p}=H_1(U_6,\mathbb F_p).
\]
It is a $37$-dimensional $\mathbb F_p$-representation of $\Gamma_6$.
To see its restriction to $C$, use the cellular chain complex of the
degree-thirty-six cover of the two-loop graph. Over a field containing
the sixth roots of unity it is
\[
 0\longrightarrow H_1(U_6)\longrightarrow k[C]^2
   \xrightarrow{\ (X-1,\,Y-1)\ }k[C]
   \longrightarrow k\longrightarrow0.
\]
Consequently the trivial character occurs twice, and each nontrivial
character of $C$ occurs exactly once. At a character $(a,b)\ne(1,1)$,
the commutator $[x,y]$ gives a nonzero vector in this one-dimensional
space: its two edge coordinates are, up to the chosen commutator sign,
$(1-b,a-1)$.

Suppose now that $p\equiv5\pmod{12}$, and let
$\omega\in\mathbb F_{p^2}$ be a primitive cube root. The two
characters
\[
 \chi=(\omega,\omega),\qquad
 \chi^{-1}=(\omega^2,\omega^2)
\]
give a two-dimensional $\mathbb F_p$-summand $V_p\subset M_p$.
It is stable under $\Gamma_6$: both characters are fixed by the
permutations of $x,y,z$, and they are interchanged by Frobenius.
Let $W_p$ be its complementary sum of character spaces, so
$M_p=V_p\oplus W_p$.

On the $\chi$ line, conjugation by $s$ acts by $-1$, while
conjugation by $r$ acts by $\omega$. Indeed
\[
 s[x,y]s^{-1}=[y,x],\qquad
 r[x,y]r^{-1}=[y,(xy)^{-1}],
\]
and the second commutator has coefficient $(XY)^{-1}$ times that
of $[x,y]$ in the metabelian commutator module. Its value at
$(\omega,\omega)$ is $\omega$. Thus the eigenvalues of $s,r,\tau$
on this line are $-1,\omega,-\omega$, respectively; none is one.
The conjugate eigenvalues occur on the $\chi^{-1}$ line.

\begin{lemma}\label{lem:cover-sign-homology}
Let $\mathcal P$ be a finite set of primes congruent to five modulo
twelve, and let
\[
 Q=\Gamma\Big/
       \bigcap_{p\in\mathcal P}\overline{U_6'U_6^p}.
\]
Then $Q$ fits into a split exact sequence
\[
 1\longrightarrow\prod_{p\in\mathcal P}M_p
   \longrightarrow Q\longrightarrow\Gamma_6\longrightarrow1,
\]
and $Q^{\mathrm{ab}}=C_6$. Choose a complement to its kernel.
For $p\in\mathcal P$, let $\alpha_p$ act as $-1$ on $V_p$, and
as the identity on $W_p$, on all $M_q$ with $q\ne p$, and on this
complement. Then $\alpha_p$ is an involutive automorphism of $Q$.
It acts trivially on $H_2(Q,\mathbb Z)$.
\end{lemma}

\begin{proof}
The simultaneous elementary abelian quotients of $U_6$ have product
image because their orders are pairwise coprime. Their kernel is the
displayed intersection, proving the exact sequence. The kernel order
is coprime to $216$, so the extension splits by the coprime splitting
theorem. The summand $V_p$ is invariant, hence its sign commutes with
the complement action and defines the asserted automorphism.
The source epimorphism gives a surjection $C_6=\Gamma^{\mathrm{ab}}
\to Q^{\mathrm{ab}}$, and the quotient to $\Gamma_6$ gives the
inverse lower bound, proving equality.

For a prime $\ell\notin\mathcal P$, the positive homology of the
kernel with $\mathbb Z_\ell$ coefficients vanishes by averaging,
so $H_2(Q,\mathbb Z)_\ell$ is identified with
$H_2(\Gamma_6,\mathbb Z)_\ell$. For $q\in\mathcal P$, remove
all other primary factors first. Since the quotient by $M_q$ has
order prime to $q$, the homology spectral sequence collapses to
\[
 H_2(Q,\mathbb Z)_q
      \simeq\bigl(H_2(M_q,\mathbb Z)\bigr)_{\Gamma_6}
      \simeq\bigl(\textstyle\bigwedge^2M_q\bigr)_{\Gamma_6}.
\]
Here the Schur multiplier of an abelian group is its exterior square;
the averaging contraction gives the spectral-sequence collapse.
The automorphism $\alpha_p$ is the identity on all terms with
$q\ne p$. For $q=p$, decompose
\[
 \bigwedge^2M_p=
   \bigwedge^2V_p\ \oplus\ (V_p\otimes W_p)\ \oplus\ \bigwedge^2W_p.
\]
It acts by $+1$ on the two exterior-square summands and by $-1$ on
the middle summand. That middle summand has zero $C$-coinvariants:
after adjoining the sixth roots, its character could be trivial
only if $W_p$ contained $\chi$ or $\chi^{-1}$, and both are in
$V_p$. Since $|C|=36$ is invertible modulo $p$, invariants and
coinvariants are exact and give the same vanishing. The middle
summand therefore contributes nothing after taking $\Gamma_6$
coinvariants. This proves the assertion about $H_2$.
\end{proof}

\subsection{Lifting the sign through central extensions}

\begin{lemma}\label{lem:sign-central-extensions}
Let $Q$ and $\alpha_p$ be as in
Lemma~\ref{lem:cover-sign-homology}, and let $q:\Gamma\to Q$ be the
natural quotient. Every finite central extension
\[
 1\longrightarrow Z\longrightarrow D\longrightarrow Q
       \longrightarrow1
\]
admits an involutive automorphism $\widetilde\alpha_p$ which lifts
$\alpha_p$ and fixes $Z$ pointwise.
For each $t\in\{s,r,\tau\}$, choose $b_t\in V_p$ such that
\[
 \alpha_p(q(t))=b_tq(t)b_t^{-1}.
\]
Such an element exists. For every lift $\widetilde t\in D$ of
$q(t)$, there is a lift $\widetilde b_t\in D$ of $b_t$ with
\[
 \widetilde\alpha_p(\widetilde t)
      =\widetilde b_t\widetilde t\widetilde b_t^{-1}.
\]
The same assertions hold with $\tau^{-1}$ as the third peripheral.
\end{lemma}

\begin{proof}
Write $Z=Z_p\times Z_{p'}$. First quotient by $Z_{p'}$.
The universal-coefficient exact sequence for trivial coefficients is
\[
 0\longrightarrow
 \operatorname{Ext}^1_{\mathbb Z}(Q^{\mathrm{ab}},Z_p)
 \longrightarrow H^2(Q,Z_p)
 \longrightarrow\operatorname{Hom}(H_2(Q,\mathbb Z),Z_p)
 \longrightarrow0.
\]
The first term is zero because $Q^{\mathrm{ab}}=C_6$ and $p\geq5$.
Lemma~\ref{lem:cover-sign-homology} therefore shows that
$\alpha_p$ fixes every class in $H^2(Q,Z_p)$. It lifts to the
given central extension with identity action on its kernel.
Such a lift is unique: two lifts differ by a homomorphism
$Q\to Z_p$, and $\operatorname{Hom}(C_6,Z_p)=0$.
Its square is consequently the identity.

For the $Z_{p'}$ extension, put $Q_0=Q/M_p$. Its order is prime to
$p$. Positive cohomology of the normal $p$-group $M_p$ with
coefficients in $Z_{p'}$ vanishes by averaging, so inflation gives
\[
 H^2(Q_0,Z_{p'})\xrightarrow{\ \sim\ }H^2(Q,Z_{p'}).
\]
The extension is therefore a pullback $Q\times_{Q_0}D_0$ of a
central extension of $Q_0$. Since $\alpha_p$ is the identity on
$Q_0$, the map $(a,d)\mapsto(\alpha_p(a),d)$ is an involutive
lift fixing $Z_{p'}$. Combining the two primary extensions over
$Q$ gives the required involution of $D$.

To construct $b_t$, use the chosen semidirect-product coordinates
of $Q$. The difference between $\alpha_p(q(t))$ and $q(t)$ lies
in $V_p$ and is minus twice the $V_p$-coordinate of $q(t)$.
Conjugation by $b\in V_p$ changes that coordinate by $(1-t)b$.
The operators $1-s,1-r,1-\tau$ are invertible on $V_p$ by the
eigenvalues computed above, so each equation has a unique solution.

Choose a lift $\widetilde b_t$ whose image in the pullback
$Q\times_{Q_0}D_0$ is $(b_t,1)$. Comparing the two lifts of
$\alpha_p(q(t))$ gives
\[
 \widetilde\alpha_p(\widetilde t)
   =\widetilde b_t\widetilde t\widetilde b_t^{-1}z
       \qquad(z\in Z).
\]
Its image in the $Z_{p'}$ extension shows that $z\in Z_p$.
Furthermore $\widetilde\alpha_p(\widetilde b_t)\widetilde b_t$
lies in the central kernel, because its image in $Q$ is $-b_t+b_t=0$.
Applying $\widetilde\alpha_p$ to the displayed equation and using
its involutivity and its identity action on $Z$ yields
$\widetilde t=\widetilde t z^2$. Squaring is injective on the
finite $p$-group $Z_p$, so $z=1$. The assertion for the inverse
peripheral follows by inversion.
\end{proof}

\begin{theorem}[Full lifts of the finite signs]
\label{thm:quadratic-sign-lifts}
For every finite set $\mathcal P$ of primes congruent to five modulo
twelve and each $p\in\mathcal P$, there is a full multiplier-one
gentle element whose action on the sixth-root homology characters
modulo each $q\in\mathcal P$ is $-1$ on
$(\omega,\omega),(\omega^2,\omega^2)$ when $q=p$, and is the
identity on all other nontrivial characters.
One may at the same time prescribe identity conjugator values in
any fixed finite quotient of $\Gamma$ of order divisible only by
two and three.
\end{theorem}

\begin{proof}
Use the finite quotient $Q$ above. Consider any finite quotient
$E$ of $\Gamma$ equipped with its quotient to $Q$, and put
$A=\ker(E\to Q)$. Lemma~\ref{lem:finite-peripheral-refinement}
constructs a finite central extension $D^\dagger\to Q$ and a
source epimorphism $q^\dagger$. Apply
Lemma~\ref{lem:sign-central-extensions} to obtain
$\widetilde\alpha_p$ and three peripheral conjugators over the
fixed $b_s,b_r,b_{\tau^{-1}}\in V_p$. Postcomposing $q^\dagger$
with this automorphism gives the new epimorphism and conjugators
required by the finite-refinement lemma. It supplies an epimorphism
to $E$ with the prescribed conjugator images in $Q$.

Now impose the peripheral product relation, generation in all finite
quotients of $\Gamma$, and the three prescribed conjugator values in
$Q$. The compactness and Hopf argument in the proof of
Theorem~\ref{thm:peripheral-lifting} applies verbatim: every finite
list is controlled by one joint quotient $E$ as just considered.
We obtain a full peripheral automorphism $\Phi$ of $\Gamma$
inducing $\alpha_p$ on $Q$.

An inner correction by the third full conjugator makes $\Phi(\tau)=\tau$.
That conjugator has image in $V_p$, so this correction acts trivially
on every $M_q$ and on $\Gamma_6$; it preserves all the homology
actions under discussion. Since the resulting automorphism is the
identity modulo $U_6$, write
\[
 \Phi(s)=sf,\qquad \Phi(r)=f^{-1}r,
       \qquad f\in U_6.
\]
The square relation says that conjugation by $s$ takes $f$ to $f^{-1}$.
Thus its image in $F^{\mathrm{ab}}$ has the form $(a,-a)$, with
$a\in6\widehat{\mathbb Z}$. Composing with conjugation by $x^a$
replaces $f$ by $y^af x^{-a}$, whose abelianization is zero.
This correction still fixes $\tau$. It acts trivially on every
$M_q$ because the deck image of $x^a$ in $C$ is the identity.
We have therefore reached $f\in F'$. The relations
$(sf)^2=(f^{-1}r)^3=1$ are exactly the normalized multiplier-one
two- and three-cycle equations. The restriction to $F$ is
\[
 x\longmapsto x,\qquad y\longmapsto f^{-1}yf,
\]
and gives the required gentle element.

For the final assertion, choose a positive integer $m$ divisible by
six and by the exponent of $T$, with no prime factors other than two
and three. Replace $U_6$ by $U_m=\overline{F'F^m}$ in the finite
construction. The deck group is now $(\mathbb Z/m\mathbb Z)^2$,
the base has order $6m^2$, and its homology still has multiplicity
one at every nontrivial character. Thus the same two cubic character
spaces define $V_p$ and all the preceding sign and homology arguments
apply without change. Take the joint image of this quotient and $T$.
Its normal kernel is still the full product of the $M_q$, because
the common quotient with $T$ kills this kernel of coprime order.
The new base is a group of order divisible only by two and three,
and its action on the homology factors through $\Gamma/U_m$.
Choose the $b_t$ in $V_p$, so their images in $T$ are the identity.
The first inner normalization has trivial image in $T$. The second
uses $a\in m\widehat{\mathbb Z}$, so $x^a$ also has trivial image
in $T$. Both therefore preserve the identity frame values.
\end{proof}

\subsection{The unitary relation without a metabelian lifting assumption}

We record the finite-cover argument which supplies the norm-one
condition. It uses only positive peripheral preservation.

\begin{lemma}\label{lem:cover-unitarity}
Let $\phi$ be a continuous automorphism of $F$ inducing the identity
on $F^{\mathrm{ab}}$ and preserving the three positive peripheral
classes. On a nontrivial sixth-root character $(a,b)$ with
\[
 a\ne1,\qquad b\ne1,\qquad ab\ne1,
\]
let $B_\phi(a,b)$ be the scalar by which it acts on the corresponding
one-dimensional summand of $H_1(U_6,\overline{\mathbb F}_p)$,
where $p\geq5$. Then
\[
 B_\phi(a,b)B_\phi(a^{-1},b^{-1})=1.
\]
\end{lemma}

\begin{proof}
The degree-thirty-six cover of the three-punctured sphere has six
boundary components over each of its three boundaries, since the
three deck elements $X,Y,XY$ have order six. It has eighteen
boundaries. Euler characteristic gives genus ten. Its profinite
fundamental group is $U_6$, and it has a surface presentation
\[
 \prod_{j=1}^{10}[a_j,b_j]\prod_{i=1}^{18}c_i=1,
\]
where the $c_i$ are positively oriented boundary elements.
Let $H=H_1(U_6,\mathbb F_p)$ and let $R$ be the span of the
boundary classes. It has dimension seventeen, their only relation
being $\sum_i[c_i]=0$.

The alternating form on $H/R$ can be seen directly in the class-two,
exponent-$p$ quotient of $U_6$. Write this quotient in logarithmic
coordinates, which are valid because $p$ is odd, and project its
degree-two commutators from $\bigwedge^2H$ to
$\bigwedge^2(H/R)$. The boundary relation gives the nondegenerate
bivector
\[
 \Omega=\sum_{j=1}^{10}[a_j]\wedge[b_j].
\]
An automorphism preserving the positive boundary conjugacy classes
permutes the $c_i$ up to conjugation. Conjugating a boundary element
changes its degree-two coordinate by a wedge involving its boundary
class, which is zero after this projection. The boundary factors'
pairwise multiplication terms are zero there for the same reason.
Applying the automorphism to the displayed surface relation therefore
gives $(\bigwedge^2\phi)\Omega=\Omega$. This is precisely preservation
of the associated nondegenerate alternating form on $H/R$.

The automorphism $\phi$ preserves $U_6$ and its positive boundary
classes, possibly permuting the components over each original
boundary. It commutes with the deck group on homology, because it
induces the identity on $F/U_6$. The boundary representation over
each original boundary is induced from the corresponding cyclic
deck subgroup. Thus a character $(a,b)$ occurs in $R$ only when
$a=1$, $b=1$, or $ab=1$; conversely each nontrivial character with
one of these equalities occurs there. The trivial representation
occurs twice in $R$. This also accounts for its dimension seventeen.

The two one-dimensional spaces at $(a,b)$ and $(a^{-1},b^{-1})$
in the statement consequently survive in $H/R$ and pair perfectly.
Deck invariance of the form pairs a character only with its inverse.
Preservation of the form forces the product of the two scalar actions
to be one, as asserted.
\end{proof}

\subsection{The quadratic characters and their independence}

For $g$ in the gentle group, use the normalized Alexander coefficient
$B_g(X,Y)$ in the convention of \cite[\S2.3, p.~101]{Furusho2003}.
At multiplier one it is the scalar just used on every nontrivial
character space. At an arbitrary multiplier $\lambda_g$, its relation
to the semilinear coefficient on the commutator is
\[
 A_g(X,Y)=S_{\lambda_g}(X)S_{\lambda_g}(Y)B_g(X,Y),
 \qquad S_a(T)=\frac{T^a-1}{T-1}.
\]
The divided power is the continuous geometric sum. Its identity
$S_{ab}(T)=S_a(T)S_b(T^a)$, together with composition of the
actions, gives
\begin{equation}
 B_{gh}(X,Y)=B_g(X,Y)
                  B_h(X^{\lambda_g},Y^{\lambda_g}).
 \label{eq:beta-crossed-law}
\end{equation}
For pentagon elements, Ihara's beta factorization gives
\begin{equation}
 B_g(a,b)B_g(ab,c)=B_g(b,c)B_g(a,bc).
 \label{eq:beta-associativity}
\end{equation}
We use here its full profinite scope, as recalled in
\cite[\S2.3, p.~101]{Furusho2003}, following
\cite[Proposition~1.6.1]{Ihara1999}.

\begin{theorem}\label{thm:quadratic-beta-family}
For every prime $p\equiv5\pmod{12}$ there is a continuous character
\[
 \vartheta_p:\widehat{\mathrm{GT}}_0\longrightarrow\mathbb F_2
\]
which vanishes on $\widehat{\mathrm{GT}}$. Their joint map is
surjective, already on the multiplier-one subgroup:
\[
 (\vartheta_p)_p:
 \widehat{\mathrm{GT}}_{0,1}
      \twoheadrightarrow\prod_{p\equiv5\ (12)}\mathbb F_2.
\]
Consequently the continuous quadratic character space of the gentle
group vanishing on the pentagon group, and also the one vanishing on
$G_{\mathbb Q}$, has dimension $\aleph_0$.
\end{theorem}

\begin{proof}
Fix $p$ and a primitive cube root $\omega\in\mathbb F_{p^2}$.
All the following values mean reduction modulo $p$ of the continuous
sixth-root evaluations. Put
\[
 E_p(g)=
 \frac{B_g(\omega,\omega)B_g(\omega^2,-1)}
      {B_g(\omega,-1)B_g(\omega,-\omega)}
        \in\mathbb F_{p^2}^{\times}.
\]
Every denominator is nonzero because these coefficients describe
invertible semilinear actions on the one-dimensional character spaces.
Equation~\eqref{eq:beta-crossed-law} gives
\[
 E_p(gh)=E_p(g)\,\sigma_{\lambda_g}(E_p(h)),
\]
where $\sigma_\lambda$ is the identity if $\lambda\equiv1\pmod3$
and Frobenius if $\lambda\equiv-1\pmod3$. Indeed every multiplier
is odd, so it fixes $-1$ and either fixes or interchanges the two
cube roots. Equation~\eqref{eq:beta-associativity}, with
$(a,b,c)=(\omega,\omega,-1)$, gives $E_p(g)=1$ for every pentagon
element.

All four character pairs in the definition satisfy the nonboundary
conditions of Lemma~\ref{lem:cover-unitarity}. For multiplier one,
that lemma says $E_p(g)E_p(g)^p=1$, since Frobenius is simultaneous
inversion on these sixth roots. For an arbitrary $g$, choose a Galois
element $h$ with the same multiplier and write $g=uh$, where
$\lambda_u=1$. The crossed law and $E_p(h)=1$ give $E_p(g)=E_p(u)$.
Thus every $E_p(g)$ lies in the norm-one subgroup, of order $p+1$.
Define
\[
 (-1)^{\vartheta_p(g)}=E_p(g)^{(p+1)/2}.
\]
This is a sign. Frobenius acts by inversion on the norm-one subgroup,
which fixes its quadratic sign, so the crossed law makes
$\vartheta_p$ an ordinary continuous homomorphism. It vanishes on
the pentagon group. Replacing $\omega$ by $\omega^2$ inverts $E_p$
and leaves the sign unchanged.

For a finite set $\mathcal P$ of these primes, the witness of
Theorem~\ref{thm:quadratic-sign-lifts} has at its selected prime $p$
the four coefficient values $(-1,1,1,1)$, and at every other prime
in $\mathcal P$ the four values $(1,1,1,1)$. Hence its $E_p$ value
is $-1$ and all the other $E_q$ values are one. Since
$p\equiv5\pmod{12}$, the integer $(p+1)/2$ is odd. The character
values on these witnesses form the identity matrix. Every finite
projection of the joint map is therefore onto; its image is compact
and hence closed, proving surjectivity onto the full product.

There are infinitely many primes congruent to five modulo twelve:
the Chebotarev theorem applied to the class $5$ in
$\operatorname{Gal}(\Q(\zeta_{12})/\Q)$ gives density $1/4$
\cite[Chapter V, Theorem 3.23, p.~164]{MilneCFT}.
The family therefore has countably infinite
rank. Conversely the marked gentle group is a compact metrizable
subspace of $\widehat{\mathbb Z}^{\times}\times\widehat F_2$.
It has countably many clopen sets, and consequently only countably
many continuous maps to $\mathbb F_2$. This proves the upper bound
for both character spaces. The Galois assertion uses
$G_{\mathbb Q}\subset\widehat{\mathrm{GT}}$.
\end{proof}

\begin{corollary}[A simultaneous profinite abelian quotient]
\label{cor:simultaneous-quotient}
There is a continuous surjection, already onto on global multiplier one,
\[
 \left(\GTz/\overline{\langle\!\langle\GT\rangle\!\rangle}\right)^{\mathrm{ab}}
 \longrightarrow
 (\Z/2\Z)^{\mathbb N}\times
 \prod_{\substack{p\geq5\\p\ \mathrm{prime}}}
       \prod_{k,j\geq1}\Z/p^k\Z.
\]
Consequently its continuous character group with values in $\Q/\Z$
contains the direct sum of the dual finite cyclic groups displayed here.
\end{corollary}
\begin{proof}
Combine Theorem~\ref{thm:quadratic-beta-family} with
Corollary~\ref{cor:all-orders}.  Every finite projection is a subdirect
product of finite primary groups of pairwise coprime orders.  Such a
subdirect product is the full product: its order is divisible by the
order of every primary factor, hence by their product.  Compactness gives
surjectivity of the infinite map.  Pullback of the finite-coordinate
characters gives the stated embedding of their direct sum.
\end{proof}
}{}
\IfFileExists{frame_character.tex}{\section{A global character from the two elliptic frames}
\label{sec:global-frame-character}

We define the frame signs, extend their sum to every multiplier, and
fix its arithmetic normalization. The construction uses the normalized
frame vanishing theorem of \cite[Theorem~3]{H17}. The extension and
the nonzero element of the next section are proved here.

\subsection{The frames and their normalization}

Let $\widehat B_3=\langle u,v\mid uvu=vuv\rangle$ be the profinite
braid group and put $c=(uv)^3$. Its central quotient is
\[
 \Gamma=\widehat B_3/\overline{\langle c\rangle}
      =\langle s,r\mid s^2=r^3=1\rangle,
 \qquad s=uvu,\quad r=uv,\quad \tau=sr=u^{-1}.
\]
The free pure subgroup $F=\ker(\Gamma\to S_3)$ has generators
$x=u^2=\tau^{-2}$ and $y=v^2=sxs$. We also use the subgroup
$\langle u^2,v^2\rangle\subset\widehat B_3$, which maps isomorphically
to $F$. A gentle element $g=(\lambda_g,f_g)$ acts on the braid
group by
\[
 \Phi_g(u)=u^{\lambda_g},\qquad
 \Phi_g(v)=f_g(u^2,v^2)^{-1}v^{\lambda_g}f_g(u^2,v^2),
 \qquad \Phi_g(c)=c^{\lambda_g}.
\]
These are automorphisms and induce the identity on the permutation
quotient \cite[\S1.3]{HS}.

The finite-subgroup and centralizer theorem for the free profinite
product gives unique pure frames
\[
 \Phi_g(s)=\beta_gs\beta_g^{-1}\quad(g\in\GTz),
 \qquad
 \Phi_h(r)=A_hrA_h^{-1}\quad(\lambda_h=1),
 \qquad \beta_g,A_h\in F.
\]
These are the frames of \cite[Lemma~5]{H17}. For the cubic frame,
the multiplier-one condition fixes the conjugacy class of $r$ in
$\Gamma^{\mathrm{ab}}=C_6$, distinguishing it from $r^{-1}$.
The frames vary continuously: their graphs are closed in compact
products, and uniqueness makes the corresponding projections
homeomorphisms. Their composition identities are
\begin{equation}
 \beta_{gh}=\Phi_g(\beta_h)\beta_g,
 \qquad A_{hk}=\Phi_h(A_k)A_h
       \quad(\lambda_h=\lambda_k=1).
 \label{eq:frame-cocycles}
\end{equation}

Let $\theta$ interchange $x,y$. Comparison of the two expressions
for $\Phi_g(y)$ gives
$f_g=\theta(\beta_g)\beta_g^{-1}$. Indeed their possible discrepancy
centralizes $y^{\lambda_g}$; abelianization makes its exponent zero.
Thus the two coordinates of $[\beta_g^{-1}]\in F^{\mathrm{ab}}$
are equal. Define $\rho_2(g)$ by
\[
 [\beta_g^{-1}]=(\rho_2(g),\rho_2(g)).
\]
It is a continuous cocycle,
\[
 \rho_2(gh)=\rho_2(g)+\lambda_g\rho_2(h).
\]
Its restriction to $G_{\mathbb Q}$ is the Kummer cocycle of the
coherent positive roots of two \cite[Lemma~7]{H17}.
Write
\[
 \mathcal G=\GTz,\qquad
 N_0=\{g\in\mathcal G:\lambda_g=1,\ \rho_2(g)=0\},
 \qquad H=G_{\mathbb Q}.
\]
The cocycle makes $N_0$ closed and normal.

We will use the arithmetic decomposition
\begin{equation}
 \mathcal G=N_0H,\qquad
 N_0\cap H=G_{K_2},\qquad
 K_2=\mathbb Q(\mu_\infty,2^{1/\infty}).
 \label{eq:gentle-arithmetic-normalization}
\end{equation}
Here is why it holds also in the gentle group. At multiplier one,
the two cycle equations make $f$ trivial in the class-two quotient
of $F$: its commutator exponent is killed by three and hence is zero
in $\widehat{\mathbb Z}$. Write the class of $\beta^{-1}$ as
$x^\rho y^\rho[x,y]^d$. The relation
$f=\theta(\beta)\beta^{-1}$ then gives
$\rho^2+2d=0$, up to the harmless common commutator sign. Thus
$\rho_2$ is even. On $G_{\mathbb Q^{\mathrm{ab}}}$ the Kummer
image of two is precisely $2\widehat{\mathbb Z}$: the positive root
$2^{1/n}$ belongs to $\mathbb Q^{\mathrm{ab}}$ exactly for $n=1,2$.
For $n>2$ its polynomial is Eisenstein, and its real root field
misses a nonreal conjugate, so is not normal and cannot be abelian.
The divisibilities of a closed subgroup of $\widehat{\mathbb Z}$
determine it, giving the asserted Kummer image.
First match the multiplier of $g$ by a Galois element; then match
the remaining even Kummer value by an element of
$G_{\mathbb Q^{\mathrm{ab}}}$. Their product gives
\eqref{eq:gentle-arithmetic-normalization}. The intersection formula
follows from the two defining Galois actions.

\subsection{The signed matrices at level twelve}

Put
\[
 U=\begin{pmatrix}1&1\\0&1\end{pmatrix},\qquad
 L=\begin{pmatrix}1&0\\-1&1\end{pmatrix},\qquad
 S=ULU,\qquad R_+=UL,\qquad R=-R_+.
\]
The assignments $u\mapsto U$, $v\mapsto L$ define the signed
representation $\widetilde q_\ell:\widehat B_3\to
\operatorname{SL}_2(\mathbb Z_\ell)$ for every prime $\ell$.
Its restriction to $F$ is denoted $q_\ell$. Reducing the signed
representation modulo twelve gives $q_{12}$ on $F$ and a projective
representation $q$ of $\Gamma$.
The signed pure image is
\[
 E_{12}=V_4\times\operatorname{SL}_2(\mathbb F_3),
 \qquad V_4=\{I,U^2,L^2,U^2L^2\}\pmod4.
\]
The matrix $-I$ is not in $V_4$, so this image embeds in the
projective group modulo twelve. Its projective centralizers of
$R_+$ and $S$ are respectively
\begin{equation}
 \{I\}\times\bigl(\langle-I\rangle\times\langle R\rangle\bigr),
 \qquad
 \{I,U^2L^2\}\times\langle S\rangle.
 \label{eq:finite-frame-centralizers}
\end{equation}
These follow directly by multiplying the indicated matrices; they
are the finite groups in \cite[Lemma~12]{H17}.

\begin{lemma}\label{lem:frame-all-multipliers-congruence}
For every $g\in\mathcal G$ and every prime $\ell$ there is a unique
matrix
\[
 D_{g,\ell}=\begin{pmatrix}\lambda_{g,\ell}&b_{g,\ell}\\0&1\end{pmatrix}
\]
such that
$\widetilde q_\ell\Phi_g=\operatorname{Ad}(D_{g,\ell})\widetilde q_\ell$.
It depends continuously on $g$, and $D_{gh,\ell}=D_{g,\ell}D_{h,\ell}$.
For $h\in N_0$, its reduction modulo twelve is the identity.
\end{lemma}

\begin{proof}
Put $a=\lambda_{g,\ell}$ and
$M=\widetilde q_\ell(\Phi_g(v))=
\left(\begin{smallmatrix}\alpha&b\\d&2-\alpha\end{smallmatrix}\right)$.
It is conjugate to $L^a$, so its trace is two and its determinant
is one. The lower-left entry of the braid relation
$U^aMU^a=MU^aM$ gives $d(1+ad)=0$.
If $d=0$, the domain property of $\mathbb Z_\ell$, together with
the trace and determinant, makes $M$ upper unipotent. The images
of both braid generators would commute, contradicting surjectivity
onto $\operatorname{SL}_2(\mathbb Z_\ell)$. Hence $d=-a^{-1}$
and $b=a(\alpha-1)^2$. Direct multiplication gives
\[
 D_{g,\ell}=
  \begin{pmatrix}a&a(1-\alpha)\\0&1\end{pmatrix}.
\]
This proves existence and continuity. A matrix commuting with both
$U,L$ is scalar; the specified lower-right entry gives uniqueness.
Composition then gives the product identity.

For $\lambda_h=1$, write the reduction modulo twelve as $U^{t_h}$.
Modulo four the conjugator $f_h\in F'$ maps trivially because
$V_4$ is abelian, so $t_h=0\pmod4$. Modulo three,
$U^{-t_h}q_3(\beta_h)$ lies in the quaternion subgroup of
$\operatorname{SL}_2(\mathbb F_3)$: its projective class centralizes
the involution represented by $S$. In the quotient of order three,
the classes of $U$ and $L$ agree, whereas
$[\beta_h]=(-\rho_2(h),-\rho_2(h))$ gives class
$-4\rho_2(h)=-\rho_2(h)$. Thus
$t_h=-\rho_2(h)\pmod3$. This is the shear formula of
\cite[Lemma~13]{H17}; in particular $D_h=I\pmod{12}$ on $N_0$.
\end{proof}

Define the closed subgroup
\[
 H_0=\{h\in\mathcal G:\lambda_h=1,\ t_h=0\pmod{12}\}.
\]
It contains $N_0$. For $h\in H_0$, the two frames lie in
\eqref{eq:finite-frame-centralizers}, so define
\begin{equation}
 q_3(A_h)\in(-I)^{a_3(h)}\langle R\rangle,
 \qquad
 q_3(\beta_h)\in S^{a_2(h)}\langle-I\rangle,
 \qquad a_2(h),a_3(h)\in\mathbb F_2.
 \label{eq:frame-sign-definition}
\end{equation}
These are exactly the signs of \cite[Definition~14]{H17}.
The cocycles \eqref{eq:frame-cocycles}, the identity action on
$E_{12}$, and the abelian groups in
\eqref{eq:finite-frame-centralizers} make $a_2,a_3$ continuous
characters on $H_0$. Put $\delta=a_2+a_3$.

The sign convention can also be recorded by a single unit. Let
$\xi(h)\in(\mathbb Z/12\mathbb Z)^\times$ be determined by
\[
 \xi(h)\equiv(-1)^{a_2(h)}\pmod4,\qquad
 \xi(h)\equiv(-1)^{a_3(h)}\pmod3.
\]
For the quadratic Dirichlet character $\chi_{12}$ the complete table is
\[
 \begin{array}{c|rrrr}
  t&1&5&7&11\\ \hline
  \chi_{12}(t)&1&-1&-1&1
 \end{array}
 \qquad (-1)^{\delta(h)}=\chi_{12}(\xi(h)).
\]
All elements of $H_0$ have multiplier one. The unit $\xi$ records
the two frame signs separately from that multiplier.

\subsection{Conjugation by elements of every multiplier}

\begin{proposition}\label{prop:frame-sign-conjugation}
The restrictions of $a_2,a_3$, and hence $\delta$, to $N_0$ are
invariant under conjugation by $\mathcal G$.
\end{proposition}

\begin{proof}
Let $h\in N_0$, $g\in\mathcal G$, and $\eta=ghg^{-1}$.
Normality gives $\eta\in N_0$. All matrix calculations in this
proof are modulo three. Write $\overline D_g$ for the reduction
of $D_{g,3}$.

Applying $\Phi_\eta\Phi_g=\Phi_g\Phi_h$ to $s$ and comparing
unique pure frames gives
\[
 \beta_\eta=\Phi_\eta(\beta_g)^{-1}\Phi_g(\beta_h)\beta_g.
\]
Since $\overline D_\eta=I$,
\[
 q_3(\beta_\eta)=q_3(\beta_g)^{-1}\overline D_g
      q_3(\beta_h)\overline D_g^{-1}q_3(\beta_g).
\]
Thus the traces of their squares agree. The defining groups give
\[
 \operatorname{tr}(q_3(\beta_h)^2)=2(-1)^{a_2(h)},
\]
which proves invariance of $a_2$.

Put $\epsilon\in\{1,-1\}$ equal to $\lambda_g$ modulo three.
The torsion theorem gives
$\Phi_g(r)=\gamma_g r^\epsilon\gamma_g^{-1}$.
The exponent follows from the action on $\Gamma^{\mathrm{ab}}=C_6$.
Since $\Phi_g$ induces the identity on $S_3$, multiplying
$\gamma_g$ by a power of $r$ permits the form
$\gamma_g=B_gs^e$, with $B_g\in F$ and $e=0$ or one.
Here $e=0$ when $\epsilon=1$ and $e=1$ when $\epsilon=-1$.
Comparison of the two actions on $r$ gives
\begin{equation}
 A_\eta=\Phi_\eta(\gamma_g)^{-1}\Phi_g(A_h)\gamma_g.
 \label{eq:cubic-frame-conjugation}
\end{equation}
The right side is pure and conjugates $r$ to $\Phi_\eta(r)$,
so equality follows from uniqueness of the pure frame.

We check this equality in the signed representation. Lift
$\gamma_g$ to $\widetilde\gamma_g=B_g(uvu)^e$ in
$\widehat B_3$. Conjugation by $uvu$ interchanges $x,y$ and
normalizes the embedded free subgroup. Since $\lambda_\eta=1$,
\[
 \Phi_\eta(uvu)=\beta_\eta(uvu)\beta_\eta^{-1}.
\]
Both sides project to the same modular element and have square $c$.
A central discrepancy would have square one, and the central
procyclic group $\overline{\langle c\rangle}\simeq\widehat{\mathbb Z}$
is torsion free. This proves the equality.
It follows that the lifted right side of
\eqref{eq:cubic-frame-conjugation} belongs to the embedded free
subgroup: the two factors involving $uvu$ cancel through its
normalizing action. Its projection equals $A_\eta$, and the free
subgroup injects into the central quotient, so it is the signed
lift of $A_\eta$ itself.

Applying the signed matrix representation therefore gives
\[
 q_3(A_\eta)=
 \widetilde q_3(\widetilde\gamma_g)^{-1}\overline D_g
 q_3(A_h)\overline D_g^{-1}
 \widetilde q_3(\widetilde\gamma_g).
\]
The trace is unchanged. Every power of $R$ has trace two in
$\mathbb F_3$, so
$\operatorname{tr}(q_3(A_h))=2(-1)^{a_3(h)}$.
This proves invariance of $a_3$ as well.
\end{proof}

\begin{theorem}[Arithmetic normalization of the frame character]
\label{thm:global-frame-character}
There is a unique continuous character
\[
 \Delta:\mathcal G\longrightarrow\mathbb F_2
\]
which restricts to $\delta$ on $N_0$ and is zero on $G_{\mathbb Q}$.
For any $g\in\mathcal G$, choose $\sigma\in G_{\mathbb Q}$
with $(\lambda,\rho_2)(\sigma)=(\lambda,\rho_2)(g)$. Then
\[
 \Delta(g)=\delta(g\sigma^{-1}).
\]
This character vanishes on the full pentagon group.
\end{theorem}

\begin{proof}
By \cite[Theorem~3]{H17}, $\delta$ vanishes on the normalized
pentagon subgroup. In particular it vanishes on
$N_0\cap H=G_{K_2}$. Using
\eqref{eq:gentle-arithmetic-normalization}, define
$\Delta(n\sigma)=\delta(n)$ for $n\in N_0$, $\sigma\in H$.
Two decompositions differ by an element of $N_0\cap H$, so this
is well defined. Proposition~\ref{prop:frame-sign-conjugation}
proves its homomorphism law. The multiplication map
$N_0\times H\to\mathcal G$ is a quotient map, since it is a
continuous surjection from a compact space to a Hausdorff space.
Thus the character is continuous. The two specified restrictions
and the decomposition prove uniqueness, as well as the displayed
formula. Finally \cite[Theorem~2]{H17} decomposes the pentagon
group into its normalized subgroup and $G_{\mathbb Q}$, both of
which have zero image. Hence $\Delta$ vanishes on $\GT$.
\end{proof}
}{}
\IfFileExists{frame_witness.tex}{\section{A frame character beyond the metabelian action}
\label{sec:frame-witness}

The quadratic beta characters are detected by metabelian homology.
The following construction gives a different kind of element: its
normalized conjugator lies in the third derived subgroup, while one
of its torsion frames has a nontrivial central value modulo twelve.
The construction uses the relative form of
Theorem~\ref{thm:peripheral-lifting}, with the closed normal subgroup
$U=F''$; no openness assumption on $U$ is needed.

Keep the modular marking $\Gamma=\langle s,r\mid s^2=r^3=1\rangle$,
$\tau=sr$, $x=\tau^{-2}$, and $y=sxs$. Let
\[
 q:\Gamma\longrightarrow
 Q_{12}=\operatorname{SL}_2(\mathbb Z/12\mathbb Z)/\{\pm I\}
\]
be the congruence quotient with the signed pure images
\[
 X=\begin{pmatrix}1&2\\0&1\end{pmatrix},\qquad
 Y=\begin{pmatrix}1&0\\-2&1\end{pmatrix}.
\]
These arise from $s=\left(\begin{smallmatrix}0&1\\-1&0\end{smallmatrix}\right)$
and $\tau=\left(\begin{smallmatrix}1&-1\\0&1\end{smallmatrix}\right)$,
with matrices taken modulo scalar sign. Put $\zeta=[5I]\in Q_{12}$.

\begin{theorem}\label{thm:deep-frame-witness}
There is a multiplier-one gentle element $e=(1,f)$ such that
\[
 f\in F''',\qquad \rho_2(e)=0.
\]
Its modular automorphism $\Phi$ fixes $\tau$ and induces the identity
under $q$. Its unique pure torsion frames, defined by
\[
 \Phi(s)=\beta s\beta^{-1},\qquad
 \Phi(r)=ArA^{-1},
\]
satisfy
\[
 q(\beta)=1,\qquad q(A)=\zeta.
\]
In particular $e$ acts identically on $F/F'''$, and every character
factoring through its metabelian action vanishes on $e$.
\end{theorem}

\begin{proof}
We first calculate the finite image needed for the lift. Set
$z=(xy)^{-1}$ and
\[
 a=[x^2,y^2],\qquad b=[y^2,z^2],\qquad w=[a,b]\in F''.
\]
Direct multiplication gives
\[
 \begin{array}{c|c|c|c}
  &q(a)&q(b)&q(w)\\ \hline
  \bmod3&
    \begin{pmatrix}1&1\\1&2\end{pmatrix}&
    \begin{pmatrix}0&2\\1&0\end{pmatrix}&-I\\[3pt]
  \bmod4&I&I&I\\[3pt]
  \bmod12&
    \begin{pmatrix}1&4\\4&5\end{pmatrix}&
    \begin{pmatrix}9&8\\4&9\end{pmatrix}&5I.
 \end{array}
\]
Modulo three, $X,Y$ generate $\operatorname{SL}_2(\mathbb F_3)$.
The two matrices in the first row generate its normal quaternion
subgroup of order eight: each squares to $-I$, they anticommute,
and conjugation by $X,Y$ preserves this subgroup. The quotient has
order three, so the derived subgroup is precisely this quaternion
group; its derived subgroup is $\{I,-I\}$. Modulo four, $X,Y$
commute and have order two. These calculations give
\[
 q(F'')\subset\{1,\zeta\},
\]
and the word $w$ gives equality. The element $\zeta$ is central
and has order two.

Put $U=F''$. For a prime $p$, let $O^p(U)$ be the intersection of
the kernels of all continuous homomorphisms from $U$ to finite
$p$-groups, and set $\Gamma_p=\Gamma/O^p(U)$.
The subgroup $O^p(U)$ is characteristic in $U$ and hence normal in
$\Gamma$. Write
\[
 R=\bigcap_{n\geq1}\gamma_n(U).
\]
There is a natural isomorphism
\begin{equation}
 \Gamma/R\simeq
   \prod_{p,\,\Gamma/U}\Gamma_p,
 \label{eq:relative-pronilpotent-product}
\end{equation}
where the right side is the fibre product with a common image in
$\Gamma/U$. Here are the details for an arbitrary closed $U$.
Its maximal pronilpotent quotient is the product of its maximal
pro-$p$ quotients: every finite subdirect product of groups of
pairwise coprime orders is the full product, and compactness passes
this surjectivity to the inverse limit. The kernel is the
intersection of the lower central series. Indeed a finite
nilpotent quotient kills some lower central term; conversely an
open subgroup containing the intersection contains some term by
compactness of a descending closed family. To obtain
\eqref{eq:relative-pronilpotent-product}, lift the common
$\Gamma/U$ coordinate to $\Gamma$ and then use the just-proved
surjectivity on $U$. The natural map is a continuous bijection
between compact Hausdorff groups and is therefore an isomorphism.

Since $q(U)=\langle\zeta\rangle$ is a two-group, the map $q$
factors through $\Gamma_2$. Let
\[
 J=\operatorname{im}(\Gamma\longrightarrow Q_{12}\times\Gamma/U).
\]
The map $\Gamma_2\to J$ has pro-two kernel. The element
$(\zeta,1)$ belongs to $J$, by the word $w$, and centralizes the
image of $r$. Choose a preimage $c_0$ in the image $U^{(2)}$ of
$U$ in $\Gamma_2$. The elements $c_0rc_0^{-1}$ and $r$ have
order three and map to the same generator of order three in $J$.
They are complements to a pro-two kernel in the preimage of that
cyclic group. Coprime complement conjugacy in every finite quotient,
followed by compactness, gives a correction in the kernel that
makes them equal. Thus there is
\[
 c\in C_{\Gamma_2}(r)\cap U^{(2)},\qquad q(c)=\zeta.
\]
For clarity, the conjugacy is by a kernel element, so it retains
the image $(\zeta,1)$; equality of the complements implies equality
of the chosen generators because both map to the same element of
order three.

In the product \eqref{eq:relative-pronilpotent-product}, let
$\overline b_r$ have coordinate $c$ at two and identity at every
other prime. It centralizes the image of $r$. Put
$\overline b_s=\overline b_{\tau^{-1}}=1$.
The identity automorphism of $\Gamma/R$ is consequently a positive
peripheral automorphism with these specified conjugators.
Theorem~\ref{thm:peripheral-lifting} supplies a full peripheral
automorphism $\Phi$ with all three specified conjugator images. Its cusp
conjugator lies in $R$; the allowed inner correction by that
conjugator makes $\Phi(\tau)=\tau$ and leaves all their images
in $\Gamma/R$, and hence in $Q_{12}$, unchanged.

The resulting conjugator of $s$ belongs to $R$, and that of $r$
belongs to $U$. Both are pure. The torsion-centralizer theorem for
a free profinite product gives $C_\Gamma(s)=\langle s\rangle$
and $C_\Gamma(r)=\langle r\rangle$, so these pure frames are
unique. They are the $\beta,A$ in the statement, with
$q(\beta)=1$ and $q(A)=\zeta$.

Since $\Phi$ is the identity modulo $R$ and fixes $\tau$, write
\[
 \Phi(s)=sf,\qquad \Phi(r)=f^{-1}r,
       \qquad f=s\Phi(s)\in R.
\]
The residual $R$ is contained in $U'$ because abelian quotients are
nilpotent. Hence $f\in U'=F'''$. The torsion relations give the
two cycle identities, and restriction to $F$ is
$x\mapsto x$, $y\mapsto f^{-1}yf$. There is no additional
abelianization correction because $f$ is already in $F'$.
This is the desired multiplier-one gentle element.

The frame definition of $\rho_2$ reads the abelianization of
$\beta^{-1}$. As $\beta\in R\subset F'''$, this abelianization
is zero, giving $\rho_2(e)=0$. The map $q$ kills $R$, so
$q\Phi=q$. Finally $f\in F'''$ gives literal identity on
$F/F'''$, which includes its metabelian quotient.
\end{proof}

\begin{corollary}\label{cor:frame-character-nonzero}
The character $\Delta$ of Theorem~\ref{thm:global-frame-character}
is nonzero on $N_0$ and is outside the span of all continuous quadratic
characters factoring through the metabelian action. In particular
$\Delta$ and the characters $\vartheta_p$ of
Theorem~\ref{thm:quadratic-beta-family} are linearly independent.
\end{corollary}

\begin{proof}
For the element of Theorem~\ref{thm:deep-frame-witness}, the two
signed matrices are $q_{12}(A)=5I$ and $q_{12}(\beta)=I$, because
the pure image $E_{12}$ embeds in the projective quotient. Hence the
frame matrices modulo three are $q_3(A)=-I$ and
$q_3(\beta)=I$. Definition~\eqref{eq:frame-sign-definition} gives
$a_3(e)=1$, $a_2(e)=0$. Thus $\xi(e)=5\pmod{12}$ and
$\Delta(e)=\delta(e)=1$. Its metabelian action is the identity,
so every character factoring through that action vanishes on it.
Evaluating a finite linear relation on $e$ first kills the coefficient
of $\Delta$; independence of the $\vartheta_p$ kills all the others.
\end{proof}
}{}

\section{Comparison with the arithmetic character question}

The abelian quotients constructed here measure the effect of omitting
the pentagon. The arithmetic question concerns a different restriction:
a character on full \(\GT\) that vanishes on its Galois subgroup.
The companion paper \cite{ArithmeticCompanion} studies that restriction
using the same cyclic-cover homology, together with CM reciprocity.
We state its conclusions with their precise scope.

Write \(G=\GT\), \(H=G_{\mathbb Q}\), and
\[
 N=\ker(\lambda,\rho_2),\qquad K=N\cap H,\qquad
 B=N/\overline{[G,N]K}.
\]
Here \(\rho_2\) is the Kummer cocycle of two in the normalization
of \cite{H17}. The decomposition \(G=NH\) gives
\[
 G^{\rm ab}\simeq\widehat{\mathbb Z}^{\times}\times B.
\]
In particular, a nonzero continuous quadratic character of \(B\)
would separate \(G\) from \(H\).
The invariant quadratic restriction
\[
 H^1(N,\mathbb F_2)^G\longrightarrow H^1(K,\mathbb F_2)^H
\]
is surjective and has an explicit linear section
\cite[Theorems~1.1 and~1.6]{ArithmeticCompanion}.
Its kernel is \(\operatorname{Hom}_{\rm cts}(B,\mathbb F_2)\),
whose vanishing remains unresolved.

The additional CM comparisons give integral arithmetic normalizations
of this same relative obstruction
\cite[Theorem~1.14]{ArithmeticCompanion}.
For \(y^p=x(1-x)\), a local reflex-norm operator determines a finite
deficit in a dyadic norm-one torus. Its determinant is governed by
the relative cyclotomic class number, but the complete arithmetic
image also depends on square classes of explicit real units.
Frobenius elements fill the deficits at \(p=29\) and \(113\).
At \(p=163\), negative-square identities prove that the arithmetic
image still has index four. Simultaneous arithmetic matching is used
only for the coordinates whose entire images are known to match.
For any collection of prime degrees, the full and arithmetic scalar
images are nevertheless exact products of their individual images.
The entire conjugation obstruction for intrinsic relative quadratic
characters vanishes, including for the infinite collection. Such a
character extends to full \(G\), zero on \(H\), precisely when its
relative central extension class is zero
\cite[Theorem~1.16]{ArithmeticCompanion}.

These results determine the arithmetic target and its lifted part
without determining the remaining relative kernel. The characters
constructed in the present paper vanish on all of \(\GT\), so their
restrictions to that kernel are zero. They prove strictness of the
gentle inclusion and provide explicit abelian quotients of its
normal quotient; an arithmetic-trivial character nonzero on full
\(\GT\) would require an additional argument.


\section{Scope and remaining questions}
Theorem~\ref{thm:weight16} computes the graded indecomposables exactly.
Proposition~\ref{prop:weight12} identifies an exact pentagon-detecting
functional.  Neither statement determines the whole profinite quotient
$\GTz/\overline{\langle\!\langle\GT\rangle\!\rangle}$, nor the
abelianization of $\GT$.  Theorem~\ref{thm:even-characters} and
Corollary~\ref{cor:all-orders}, together with
Theorem~\ref{thm:quadratic-beta-family}, construct large continuous quotients after
the pentagon group is killed; their targets do not exhaust the whole
quotient.  A character vanishing on $\GT$ cannot separate
$\GT$ from its Galois subgroup.  The arithmetic question asks instead
for a character that vanishes on the Galois subgroup and is nonzero
on $\GT$.
The prime-three character group and the full nonabelian structure of
the gentle quotient also remain undetermined here.

\section*{Note on method}
The mathematics and text were produced by GPT-6 Astra Pro (ChatGPT Pro mode).
The finite Lie-space and pentagon calculations were independently
implemented in this work; the precise proof-check events are listed at
the \href{\recordurl}{permanent records address}.
The frame-character construction uses the normalized vanishing theorem
of \cite{H17}; its extension to all multipliers and its nonzero full
profinite witness are proved in Sections~\ref{sec:global-frame-character}
and~\ref{sec:frame-witness}.
The quadratic finite-cover construction and the complete global frame
theorem received separate statement-only reconstructions by GPT-6 Astra
Pro, followed by full targeted comparisons. The latter check includes
every multiplier, the unique arithmetic extension, and the witness in
the third derived subgroup. The effort setting is not exposed. The full peripheral lifting theorem
received a complete fresh statement-only reconstruction and a full
targeted comparison. The new proof uses finite central character
extensions, counting, and compactness. Its initial
incomplete attempt remains identified as partial in the records. The cited monoid theorem records
a proof by Claude Fable 5.2 and a check by Claude Opus 5.5, both at
effort max. The full odd-primary product theorem, integral top-weight
transfer for every prime at least five, and dual-space statement also
received complete fresh statement-only reconstructions and full
comparisons. The complete weight-sixteen and weight-twelve statements
also received fresh independent reconstructions with separate code,
explicit minors and exact integer identities.
The cited H17 note records the same two models at effort
max for its stated background checks.
No human has checked this paper.
The text is dedicated to the public domain under CC0 1.0.
The associated \href{\workbenchurl}{workbench address} is designated
for the program \emph{Gentle derivations and pentagon certificates},
its exact inputs and outputs, and the records.

\begin{thebibliography}{99}
\bibitem{CohenFalkRandell} D.~C.~Cohen, M.~Falk and R.~Randell,
Pure braid groups are not residually free, in \emph{Configuration Spaces},
Edizioni della Normale (2012), 213--230, \S4, p.~223;
\href{https://www.math.lsu.edu/~cohen/dcpapers/CohenFalkRandellResFree.pdf}{author's text}.
\bibitem{H17} \emph{Relations $(\mathrm{III}')$ and $(\mathrm{IV})$ of
Nakamura and Schneps hold on the profinite Grothendieck--Teichm\"uller
group, with an application to the two elliptic points of the modular
curve}, version 2, 24 November 2026,
\href{https://hexagonmath.org/2611.00940v2}{hexagon:2611.00940v2};
Theorems 2--3, Lemmas 5, 7, 12--13, and Definition 14.
\bibitem{Monoid} \emph{The profinite Grothendieck--Teichm\"uller monoid is a group, on a question of Lochak and Schneps},
25 November 2026, \href{https://hexagonmath.org/2611.01024}{hexagon:2611.01024},
Theorems A and E; proof by Claude Fable 5.2 and independent check by
Claude Opus 5.5, effort max, no human verification.
\bibitem{Drinfeld} V.~G.~Drinfeld, On quasitriangular quasi-Hopf algebras
and on a group that is closely connected with $\operatorname{Gal}(\overline\Q/\Q)$,
Leningrad Math. J. \textbf{2} (1991), 829--860; \S5.
\bibitem{HS} D.~Harbater and L.~Schneps, Fundamental groups of moduli
and the Grothendieck--Teichm\"uller group, Trans. Amer. Math. Soc.
\textbf{352} (2000), 3117--3148; Main Theorem and \S1.3.
\bibitem{DG} V.~A.~Dolgushev and J.~J.~Guynee, GT-shadows for the gentle
version of the Grothendieck--Teichmueller group,
J. Pure Appl. Algebra \textbf{229} (2025), 107819;
\href{https://arxiv.org/abs/2401.06870}{arXiv:2401.06870v2},
Definition 2.5, Theorem 5.2 and Corollary 5.4.
\bibitem{Furusho} H.~Furusho, Pentagon and hexagon equations,
Ann. of Math. \textbf{171} (2010), 545--556;
\href{https://arxiv.org/abs/math/0702128}{arXiv:math/0702128v3}.
\bibitem{Furusho2003} H.~Furusho, Geometric and arithmetic subgroups
of the Grothendieck--Teichm\"uller group, Math. Res. Lett. \textbf{10}
(2003), 97--108, \S2.3, p.~101;
\href{https://www.math.nagoya-u.ac.jp/~furusho/paper/GTA-mrl.pdf}{author's published text}.
\bibitem{Ihara1999} Y.~Ihara, On beta and gamma functions associated
with the Grothendieck--Teichm\"uller group, in \emph{Aspects of Galois
Theory}, London Math. Soc. Lecture Note Ser. 256, Cambridge University
Press, 1999, 144--179, Proposition 1.6.1.
\bibitem{MilneCFT} J.~S.~Milne, \emph{Class Field Theory}, version 4.03
(2020), Chapter V, Theorem 3.23, p.~164;
\href{https://www.jmilne.org/math/CourseNotes/CFT.pdf}{author's course notes}.
\bibitem{Brown} F.~Brown, Mixed Tate motives over $\Z$,
Ann. of Math. \textbf{175} (2012), 949--976;
\href{https://arxiv.org/abs/1102.1312}{arXiv:1102.1312}, Theorem 1.1 and \S1.1.
\bibitem{Willwacher} T.~Willwacher, The Deligne--Drinfeld conjecture
after OpenAI, \href{https://arxiv.org/abs/2610.15592}{arXiv:2610.15592v4}
(28 October 2026), Theorem 1.1 and equation (7), pp.~1--2.
\bibitem{LazardTransport} O.~Garaialde Oca\~na and J.~Gonz\'alez-S\'anchez,
Transporting cohomology in Lazard correspondence,
\href{https://arxiv.org/abs/1405.4654}{arXiv:1405.4654v2},
Theorem 1, p.~3.
\bibitem{Markarian} N.~Markarian, A remark on the AI proof of the
Deligne--Drinfeld conjecture,
\href{https://arxiv.org/abs/2611.01745}{arXiv:2611.01745v2}
(18 November 2026), Theorem 4.8, \S5 and Appendix C.
\bibitem{OpenAI} OpenAI, The Deligne--Drinfeld conjecture, 23 September
2026, Theorem 1.1, p.~2; \href{https://github.com/openai/math/blob/main/preprints/The-Deligne-Drinfeld-conjecture-September-23-2026/paper.pdf}{OpenAI Math Release}.
\bibitem{ArithmeticCompanion} \emph{The arithmetic quotient of the profinite
Grothendieck--Teichm\"uller group}, companion paper, 25 November 2026,
Theorems~1.1, 1.6 and~1.14, \url{\workbenchurl}.
\end{thebibliography}
% Generated after the source commit; excluded from the source manifest.
\par\bigskip\begingroup\footnotesize
\noindent\textbf{Version identifiers.} UNNUMBERED-GC/v1. The following SHA-256 identifies the source manifest, whose scope excludes the PDF and generated identifiers. The complete payload manifest is available separately.\par
\noindent Source manifest SHA-256:\par\noindent\texttt{56d6feb293a11c18\allowbreak{}579a18d3589bfac6\allowbreak{}4f256327f6cfa449\allowbreak{}f709127faeb23d80}.\par
\noindent Local source commit:\par\noindent\texttt{d6e646cedccfd6aa\allowbreak{}897b9acc2fbbdfd9\allowbreak{}355c3716}.\par
This commit precedes generation of these identifiers and the PDF; the final release commit is recorded by the version tag.\par\endgroup

\end{document}
